% Môn: Vật lý
% Lớp: 10
% Chương: Chương III. Động lực học
% Bài: L10C3 Bài 20. Một số ví dụ về cách giải các bài toán thuộc phần động lực học
% Số câu: 162
% App đọc file này trực tiếp — sửa rồi Commit trên GitHub, bấm Đồng bộ GitHub .tex

% L10C3 Bài 20. Một số ví dụ về cách giải các bài toán thuộc phần động lực học

% ===== Câu 1 =====
% ID: AIQ_L10C3_FULL_1135
% Mức: NB
\dangbt{Phân tích lực và lập phương trình định luật II Newton}
\begin{ex}
% ID: AIQ_L10C3_FULL_1135
% Mức: NB
Vật khối lượng $3\,\text{kg}$ chịu hợp lực $19\,\text{N}$. Tính gia tốc.
\choice
{$8,33 m/s^2$}
{\True $6,33 m/s^2$}
{$7,6 m/s^2$}
{$5,06 m/s^2$}
\loigiai{
Theo định luật II Newton: $a=F/m=19/3=6,33$ m/s^2. Chọn B.
}
\end{ex}

% ===== Câu 2 =====
% ID: AIQ_L10C3_FULL_1136
% Mức: NB
\begin{ex}
% ID: AIQ_L10C3_FULL_1136
% Mức: NB
Vật khối lượng $4\,\text{kg}$ chịu hợp lực $13\,\text{N}$. Tính gia tốc.
\choice
{\True $3,25 m/s^2$}
{$3,9 m/s^2$}
{$2,6 m/s^2$}
{$5,25 m/s^2$}
\loigiai{
Theo định luật II Newton: $a=F/m=13/4=3,25$ m/s^2. Chọn A.
}
\end{ex}

% ===== Câu 3 =====
% ID: AIQ_L10C3_FULL_1137
% Mức: NB
\begin{ex}
% ID: AIQ_L10C3_FULL_1137
% Mức: NB
Vật khối lượng $5\,\text{kg}$ chịu hợp lực $15\,\text{N}$. Tính gia tốc.
\choice
{$3,6 m/s^2$}
{$2,4 m/s^2$}
{$5 m/s^2$}
{\True $3 m/s^2$}
\loigiai{
Theo định luật II Newton: $a=F/m=15/5=3$ m/s^2. Chọn D.
}
\end{ex}

% ===== Câu 4 =====
% ID: AIQ_L10C3_FULL_1138
% Mức: TH
\begin{ex}
% ID: AIQ_L10C3_FULL_1138
% Mức: TH
Vật khối lượng $6\,\text{kg}$ chịu hợp lực $17\,\text{N}$. Tính gia tốc.
\choice
{$2,26 m/s^2$}
{$4,83 m/s^2$}
{\True $2,83 m/s^2$}
{$3,4 m/s^2$}
\loigiai{
Theo định luật II Newton: $a=F/m=17/6=2,83$ m/s^2. Chọn C.
}
\end{ex}

% ===== Câu 5 =====
% ID: AIQ_L10C3_FULL_1139
% Mức: TH
\begin{ex}
% ID: AIQ_L10C3_FULL_1139
% Mức: TH
Vật khối lượng $7\,\text{kg}$ chịu hợp lực $19\,\text{N}$. Tính gia tốc.
\choice
{$4,71 m/s^2$}
{\True $2,71 m/s^2$}
{$3,25 m/s^2$}
{$2,17 m/s^2$}
\loigiai{
Theo định luật II Newton: $a=F/m=19/7=2,71$ m/s^2. Chọn B.
}
\end{ex}

% ===== Câu 6 =====
% ID: AIQ_L10C3_FULL_1140
% Mức: TH
\begin{ex}
% ID: AIQ_L10C3_FULL_1140
% Mức: TH
Vật khối lượng $8\,\text{kg}$ chịu hợp lực $13\,\text{N}$. Tính gia tốc.
\choice
{\True $1,63 m/s^2$}
{$1,96 m/s^2$}
{$1,3 m/s^2$}
{$3,63 m/s^2$}
\loigiai{
Theo định luật II Newton: $a=F/m=13/8=1,63$ m/s^2. Chọn A.
}
\end{ex}

% ===== Câu 7 =====
% ID: AIQ_L10C3_FULL_1141
% Mức: TH
\begin{ex}
% ID: AIQ_L10C3_FULL_1141
% Mức: TH
Vật khối lượng $2\,\text{kg}$ chịu hợp lực $15\,\text{N}$. Tính gia tốc.
\choice
{$9 m/s^2$}
{$6 m/s^2$}
{$9,5 m/s^2$}
{\True $7,5 m/s^2$}
\loigiai{
Theo định luật II Newton: $a=F/m=15/2=7,5$ m/s^2. Chọn D.
}
\end{ex}

% ===== Câu 8 =====
% ID: AIQ_L10C3_FULL_1142
% Mức: VD
\begin{ex}
% ID: AIQ_L10C3_FULL_1142
% Mức: VD
Vật khối lượng $3\,\text{kg}$ chịu hợp lực $17\,\text{N}$. Tính gia tốc.
\choice
{$4,54 m/s^2$}
{$7,67 m/s^2$}
{\True $5,67 m/s^2$}
{$6,8 m/s^2$}
\loigiai{
Theo định luật II Newton: $a=F/m=17/3=5,67$ m/s^2. Chọn C.
}
\end{ex}

% ===== Câu 9 =====
% ID: AIQ_L10C3_FULL_1143
% Mức: VD
\begin{ex}
% ID: AIQ_L10C3_FULL_1143
% Mức: VD
Vật khối lượng $4\,\text{kg}$ chịu hợp lực $19\,\text{N}$. Tính gia tốc.
\choice
{$6,75 m/s^2$}
{\True $4,75 m/s^2$}
{$5,7 m/s^2$}
{$3,8 m/s^2$}
\loigiai{
Theo định luật II Newton: $a=F/m=19/4=4,75$ m/s^2. Chọn B.
}
\end{ex}

% ===== Câu 10 =====
% ID: AIQ_L10C3_FULL_1144
% Mức: VDC
\begin{ex}
% ID: AIQ_L10C3_FULL_1144
% Mức: VDC
Vật khối lượng $5\,\text{kg}$ chịu hợp lực $13\,\text{N}$. Tính gia tốc.
\choice
{\True $2,6 m/s^2$}
{$3,12 m/s^2$}
{$2,08 m/s^2$}
{$4,6 m/s^2$}
\loigiai{
Theo định luật II Newton: $a=F/m=13/5=2,6$ m/s^2. Chọn A.
}
\end{ex}

% ===== Câu 11 =====
% ID: AIQ_L10C3_FULL_1145
% Mức: NB
\begin{ex}
% ID: AIQ_L10C3_FULL_1145
% Mức: NB
Vật khối lượng $6\,\text{kg}$ chịu hợp lực $15\,\text{N}$. Tính gia tốc.
\choiceTF
{\True Kết quả đúng là $2,5\,\text{m}$ /s^2.}
{Kết quả bằng $5\,\text{m}$ /s^2.}
{\True Khi hợp lực bằng 0, vật có gia tốc bằng 0.}
{Có thể bỏ qua đơn vị của đại lượng trong kết quả vật lí.}
\loigiai{
\begin{itemchoice}
\item (Đúng) Kết quả đúng là $2,5\,\text{m}$ /s^2. (Vì): Theo định luật II Newton: $a=F/m=15/6=2,5$ m/s^2. b) Sai vì giá trị hoặc chiều nêu ra không phù hợp. c) Đúng theo định luật II Newton. d) Sai vì kết quả vật lí phải kèm đơn vị thích hợp
\item (Sai) Kết quả bằng $5\,\text{m}$ /s^2.
\item (Đúng) Khi hợp lực bằng 0, vật có gia tốc bằng 0.
\item (Sai) Có thể bỏ qua đơn vị của đại lượng trong kết quả vật lí.
\end{itemchoice}
}
\end{ex}

% ===== Câu 12 =====
% ID: AIQ_L10C3_FULL_1146
% Mức: NB
\begin{ex}
% ID: AIQ_L10C3_FULL_1146
% Mức: NB
Vật khối lượng $7\,\text{kg}$ chịu hợp lực $17\,\text{N}$. Tính gia tốc.
\choiceTF
{\True Kết quả đúng là $2,43\,\text{m}$ /s^2.}
{Kết quả bằng $4,86\,\text{m}$ /s^2.}
{\True Khi hợp lực bằng 0, vật có gia tốc bằng 0.}
{Có thể bỏ qua đơn vị của đại lượng trong kết quả vật lí.}
\loigiai{
\begin{itemchoice}
\item (Đúng) Kết quả đúng là $2,43\,\text{m}$ /s^2. (Vì): Theo định luật II Newton: $a=F/m=17/7=2,43$ m/s^2. b) Sai vì giá trị hoặc chiều nêu ra không phù hợp. c) Đúng theo định luật II Newton. d) Sai vì kết quả vật lí phải kèm đơn vị thích hợp
\item (Sai) Kết quả bằng $4,86\,\text{m}$ /s^2.
\item (Đúng) Khi hợp lực bằng 0, vật có gia tốc bằng 0.
\item (Sai) Có thể bỏ qua đơn vị của đại lượng trong kết quả vật lí.
\end{itemchoice}
}
\end{ex}

% ===== Câu 13 =====
% ID: AIQ_L10C3_FULL_1147
% Mức: TH
\begin{ex}
% ID: AIQ_L10C3_FULL_1147
% Mức: TH
Vật khối lượng $8\,\text{kg}$ chịu hợp lực $19\,\text{N}$. Tính gia tốc.
\choiceTF
{\True Kết quả đúng là $2,38\,\text{m}$ /s^2.}
{Kết quả bằng $4,76\,\text{m}$ /s^2.}
{\True Khi hợp lực bằng 0, vật có gia tốc bằng 0.}
{Có thể bỏ qua đơn vị của đại lượng trong kết quả vật lí.}
\loigiai{
\begin{itemchoice}
\item (Đúng) Kết quả đúng là $2,38\,\text{m}$ /s^2. (Vì): Theo định luật II Newton: $a=F/m=19/8=2,38$ m/s^2. b) Sai vì giá trị hoặc chiều nêu ra không phù hợp. c) Đúng theo định luật II Newton. d) Sai vì kết quả vật lí phải kèm đơn vị thích hợp
\item (Sai) Kết quả bằng $4,76\,\text{m}$ /s^2.
\item (Đúng) Khi hợp lực bằng 0, vật có gia tốc bằng 0.
\item (Sai) Có thể bỏ qua đơn vị của đại lượng trong kết quả vật lí.
\end{itemchoice}
}
\end{ex}

% ===== Câu 14 =====
% ID: AIQ_L10C3_FULL_1148
% Mức: TH
\begin{ex}
% ID: AIQ_L10C3_FULL_1148
% Mức: TH
Vật khối lượng $2\,\text{kg}$ chịu hợp lực $13\,\text{N}$. Tính gia tốc.
\choiceTF
{\True Kết quả đúng là $6,5\,\text{m}$ /s^2.}
{Kết quả bằng $13\,\text{m}$ /s^2.}
{\True Khi hợp lực bằng 0, vật có gia tốc bằng 0.}
{Có thể bỏ qua đơn vị của đại lượng trong kết quả vật lí.}
\loigiai{
\begin{itemchoice}
\item (Đúng) Kết quả đúng là $6,5\,\text{m}$ /s^2. (Vì): Theo định luật II Newton: $a=F/m=13/2=6,5$ m/s^2. b) Sai vì giá trị hoặc chiều nêu ra không phù hợp. c) Đúng theo định luật II Newton. d) Sai vì kết quả vật lí phải kèm đơn vị thích hợp
\item (Sai) Kết quả bằng $13\,\text{m}$ /s^2.
\item (Đúng) Khi hợp lực bằng 0, vật có gia tốc bằng 0.
\item (Sai) Có thể bỏ qua đơn vị của đại lượng trong kết quả vật lí.
\end{itemchoice}
}
\end{ex}

% ===== Câu 15 =====
% ID: AIQ_L10C3_FULL_1149
% Mức: VD
\begin{ex}
% ID: AIQ_L10C3_FULL_1149
% Mức: VD
Vật khối lượng $3\,\text{kg}$ chịu hợp lực $15\,\text{N}$. Tính gia tốc.
\choiceTF
{\True Kết quả đúng là $5\,\text{m}$ /s^2.}
{Kết quả bằng $10\,\text{m}$ /s^2.}
{\True Khi hợp lực bằng 0, vật có gia tốc bằng 0.}
{Có thể bỏ qua đơn vị của đại lượng trong kết quả vật lí.}
\loigiai{
\begin{itemchoice}
\item (Đúng) Kết quả đúng là $5\,\text{m}$ /s^2. (Vì): Theo định luật II Newton: $a=F/m=15/3=5$ m/s^2. b) Sai vì giá trị hoặc chiều nêu ra không phù hợp. c) Đúng theo định luật II Newton. d) Sai vì kết quả vật lí phải kèm đơn vị thích hợp
\item (Sai) Kết quả bằng $10\,\text{m}$ /s^2.
\item (Đúng) Khi hợp lực bằng 0, vật có gia tốc bằng 0.
\item (Sai) Có thể bỏ qua đơn vị của đại lượng trong kết quả vật lí.
\end{itemchoice}
}
\end{ex}

% ===== Câu 16 =====
% ID: AIQ_L10C3_FULL_1150
% Mức: TH
\begin{ex}
% ID: AIQ_L10C3_FULL_1150
% Mức: TH
Vật khối lượng $4\,\text{kg}$ chịu hợp lực $17\,\text{N}$. Tính gia tốc.
\shortans{4,25}
\loigiai{
Theo định luật II Newton: $a=F/m=17/4=4,25$ m/s^2. Đáp án trả lời ngắn không ghi đơn vị.
}
\end{ex}

% ===== Câu 17 =====
% ID: AIQ_L10C3_FULL_1151
% Mức: TH
\begin{ex}
% ID: AIQ_L10C3_FULL_1151
% Mức: TH
Vật khối lượng $5\,\text{kg}$ chịu hợp lực $19\,\text{N}$. Tính gia tốc.
\shortans{3,8}
\loigiai{
Theo định luật II Newton: $a=F/m=19/5=3,8$ m/s^2. Đáp án trả lời ngắn không ghi đơn vị.
}
\end{ex}

% ===== Câu 18 =====
% ID: AIQ_L10C3_FULL_1152
% Mức: VD
\begin{ex}
% ID: AIQ_L10C3_FULL_1152
% Mức: VD
Vật khối lượng $6\,\text{kg}$ chịu hợp lực $13\,\text{N}$. Tính gia tốc.
\shortans{2,17}
\loigiai{
Theo định luật II Newton: $a=F/m=13/6=2,17$ m/s^2. Đáp án trả lời ngắn không ghi đơn vị.
}
\end{ex}

% ===== Câu 19 =====
% ID: AIQ_L10C3_FULL_1153
% Mức: VD
\begin{ex}
% ID: AIQ_L10C3_FULL_1153
% Mức: VD
Vật khối lượng $7\,\text{kg}$ chịu hợp lực $15\,\text{N}$. Tính gia tốc.
\shortans{2,14}
\loigiai{
Theo định luật II Newton: $a=F/m=15/7=2,14$ m/s^2. Đáp án trả lời ngắn không ghi đơn vị.
}
\end{ex}

% ===== Câu 20 =====
% ID: AIQ_L10C3_FULL_1154
% Mức: VD
\begin{ex}
% ID: AIQ_L10C3_FULL_1154
% Mức: VD
Vật khối lượng $8\,\text{kg}$ chịu hợp lực $17\,\text{N}$. Tính gia tốc.
\shortans{2,13}
\loigiai{
Theo định luật II Newton: $a=F/m=17/8=2,13$ m/s^2. Đáp án trả lời ngắn không ghi đơn vị.
}
\end{ex}

% ===== Câu 21 =====
% ID: AIQ_L10C3_FULL_1155
% Mức: VDC
\begin{ex}
% ID: AIQ_L10C3_FULL_1155
% Mức: VDC
Vật khối lượng $2\,\text{kg}$ chịu hợp lực $19\,\text{N}$. Tính gia tốc.
\shortans{9,5}
\loigiai{
Theo định luật II Newton: $a=F/m=19/2=9,5$ m/s^2. Đáp án trả lời ngắn không ghi đơn vị.
}
\end{ex}

% ===== Câu 22 =====
% ID: AIQ_L10C3_FULL_1156
% Mức: TH
\begin{bt}
% ID: AIQ_L10C3_FULL_1156
% Mức: TH
Vật khối lượng $3\,\text{kg}$ chịu hợp lực $13\,\text{N}$. Tính gia tốc.
\loigiai{
Phân tích dữ kiện, chọn định luật phù hợp. Theo định luật II Newton: $a=F/m=13/3=4,33$ m/s^2.
}
\end{bt}

% ===== Câu 23 =====
% ID: AIQ_L10C3_FULL_1157
% Mức: TH
\begin{bt}
% ID: AIQ_L10C3_FULL_1157
% Mức: TH
Vật khối lượng $4\,\text{kg}$ chịu hợp lực $15\,\text{N}$. Tính gia tốc.
\loigiai{
Phân tích dữ kiện, chọn định luật phù hợp. Theo định luật II Newton: $a=F/m=15/4=3,75$ m/s^2.
}
\end{bt}

% ===== Câu 24 =====
% ID: AIQ_L10C3_FULL_1158
% Mức: VD
\begin{bt}
% ID: AIQ_L10C3_FULL_1158
% Mức: VD
Vật khối lượng $5\,\text{kg}$ chịu hợp lực $17\,\text{N}$. Tính gia tốc.
\loigiai{
Phân tích dữ kiện, chọn định luật phù hợp. Theo định luật II Newton: $a=F/m=17/5=3,4$ m/s^2.
}
\end{bt}

% ===== Câu 25 =====
% ID: AIQ_L10C3_FULL_1159
% Mức: VD
\begin{bt}
% ID: AIQ_L10C3_FULL_1159
% Mức: VD
Vật khối lượng $6\,\text{kg}$ chịu hợp lực $19\,\text{N}$. Tính gia tốc.
\loigiai{
Phân tích dữ kiện, chọn định luật phù hợp. Theo định luật II Newton: $a=F/m=19/6=3,17$ m/s^2.
}
\end{bt}

% ===== Câu 26 =====
% ID: AIQ_L10C3_FULL_1160
% Mức: VD
\begin{bt}
% ID: AIQ_L10C3_FULL_1160
% Mức: VD
Vật khối lượng $7\,\text{kg}$ chịu hợp lực $13\,\text{N}$. Tính gia tốc.
\loigiai{
Phân tích dữ kiện, chọn định luật phù hợp. Theo định luật II Newton: $a=F/m=13/7=1,86$ m/s^2.
}
\end{bt}

% ===== Câu 27 =====
% ID: AIQ_L10C3_FULL_1161
% Mức: VDC
\begin{bt}
% ID: AIQ_L10C3_FULL_1161
% Mức: VDC
Vật khối lượng $8\,\text{kg}$ chịu hợp lực $15\,\text{N}$. Tính gia tốc.
\loigiai{
Phân tích dữ kiện, chọn định luật phù hợp. Theo định luật II Newton: $a=F/m=15/8=1,88$ m/s^2.
}
\end{bt}

% ===== Câu 28 =====
% ID: AIQ_L10C3_FULL_1162
% Mức: NB
\dangbt{Chuyển động của vật trên mặt phẳng ngang có ma sát}
\begin{ex}
% ID: AIQ_L10C3_FULL_1162
% Mức: NB
Kéo ngang vật $2\,\text{kg}$ bằng lực 27 $N$ trên sàn có hệ số ma sát 0,25. Lấy g= $10\,\text{m}$ /s^2. Tính gia tốc.
\choice
{$8,8 m/s^2$}
{$13 m/s^2$}
{\True $11 m/s^2$}
{$13,2 m/s^2$}
\loigiai{
$f=\mu mg=5N$; $a=(F-f)/m=11$ m/s^2. Chọn C.
}
\end{ex}

% ===== Câu 29 =====
% ID: AIQ_L10C3_FULL_1163
% Mức: NB
\begin{ex}
% ID: AIQ_L10C3_FULL_1163
% Mức: NB
Kéo ngang vật $3\,\text{kg}$ bằng lực 29 $N$ trên sàn có hệ số ma sát 0,3. Lấy g= $10\,\text{m}$ /s^2. Tính gia tốc.
\choice
{$8,67 m/s^2$}
{\True $6,67 m/s^2$}
{$8 m/s^2$}
{$5,34 m/s^2$}
\loigiai{
$f=\mu mg=9N$; $a=(F-f)/m=6,67$ m/s^2. Chọn B.
}
\end{ex}

% ===== Câu 30 =====
% ID: AIQ_L10C3_FULL_1164
% Mức: NB
\begin{ex}
% ID: AIQ_L10C3_FULL_1164
% Mức: NB
Kéo ngang vật $4\,\text{kg}$ bằng lực 23 $N$ trên sàn có hệ số ma sát 0,4. Lấy g= $10\,\text{m}$ /s^2. Tính gia tốc.
\choice
{\True $1,75 m/s^2$}
{$2,1 m/s^2$}
{$1,4 m/s^2$}
{$3,75 m/s^2$}
\loigiai{
$f=\mu mg=16N$; $a=(F-f)/m=1,75$ m/s^2. Chọn A.
}
\end{ex}

% ===== Câu 31 =====
% ID: AIQ_L10C3_FULL_1165
% Mức: TH
\begin{ex}
% ID: AIQ_L10C3_FULL_1165
% Mức: TH
Kéo ngang vật $5\,\text{kg}$ bằng lực 25 $N$ trên sàn có hệ số ma sát 0,1. Lấy g= $10\,\text{m}$ /s^2. Tính gia tốc.
\choice
{$4,8 m/s^2$}
{$3,2 m/s^2$}
{$6 m/s^2$}
{\True $4 m/s^2$}
\loigiai{
$f=\mu mg=5N$; $a=(F-f)/m=4$ m/s^2. Chọn D.
}
\end{ex}

% ===== Câu 32 =====
% ID: AIQ_L10C3_FULL_1166
% Mức: TH
\begin{ex}
% ID: AIQ_L10C3_FULL_1166
% Mức: TH
Kéo ngang vật $6\,\text{kg}$ bằng lực 27 $N$ trên sàn có hệ số ma sát 0,2. Lấy g= $10\,\text{m}$ /s^2. Tính gia tốc.
\choice
{$2 m/s^2$}
{$4,5 m/s^2$}
{\True $2,5 m/s^2$}
{$3 m/s^2$}
\loigiai{
$f=\mu mg=12N$; $a=(F-f)/m=2,5$ m/s^2. Chọn C.
}
\end{ex}

% ===== Câu 33 =====
% ID: AIQ_L10C3_FULL_1167
% Mức: TH
\begin{ex}
% ID: AIQ_L10C3_FULL_1167
% Mức: TH
Kéo ngang vật $7\,\text{kg}$ bằng lực 29 $N$ trên sàn có hệ số ma sát 0,25. Lấy g= $10\,\text{m}$ /s^2. Tính gia tốc.
\choice
{$3,64 m/s^2$}
{\True $1,64 m/s^2$}
{$1,97 m/s^2$}
{$1,31 m/s^2$}
\loigiai{
$f=\mu mg=17,5N$; $a=(F-f)/m=1,64$ m/s^2. Chọn B.
}
\end{ex}

% ===== Câu 34 =====
% ID: AIQ_L10C3_FULL_1168
% Mức: TH
\begin{ex}
% ID: AIQ_L10C3_FULL_1168
% Mức: TH
Kéo ngang vật $8\,\text{kg}$ bằng lực 23 $N$ trên sàn có hệ số ma sát 0,3. Lấy g= $10\,\text{m}$ /s^2. Tính gia tốc.
\choice
{\True $-0,12 m/s^2$}
{$-0,14 m/s^2$}
{$0 m/s^2$}
{$1,88 m/s^2$}
\loigiai{
$f=\mu mg=24N$; $a=(F-f)/m=-0,12$ m/s^2. Chọn A.
}
\end{ex}

% ===== Câu 35 =====
% ID: AIQ_L10C3_FULL_1169
% Mức: VD
\begin{ex}
% ID: AIQ_L10C3_FULL_1169
% Mức: VD
Kéo ngang vật $2\,\text{kg}$ bằng lực 25 $N$ trên sàn có hệ số ma sát 0,4. Lấy g= $10\,\text{m}$ /s^2. Tính gia tốc.
\choice
{$10,2 m/s^2$}
{$6,8 m/s^2$}
{$10,5 m/s^2$}
{\True $8,5 m/s^2$}
\loigiai{
$f=\mu mg=8N$; $a=(F-f)/m=8,5$ m/s^2. Chọn D.
}
\end{ex}

% ===== Câu 36 =====
% ID: AIQ_L10C3_FULL_1170
% Mức: VD
\begin{ex}
% ID: AIQ_L10C3_FULL_1170
% Mức: VD
Kéo ngang vật $3\,\text{kg}$ bằng lực 27 $N$ trên sàn có hệ số ma sát 0,1. Lấy g= $10\,\text{m}$ /s^2. Tính gia tốc.
\choice
{$6,4 m/s^2$}
{$10 m/s^2$}
{\True $8 m/s^2$}
{$9,6 m/s^2$}
\loigiai{
$f=\mu mg=3N$; $a=(F-f)/m=8$ m/s^2. Chọn C.
}
\end{ex}

% ===== Câu 37 =====
% ID: AIQ_L10C3_FULL_1171
% Mức: VDC
\begin{ex}
% ID: AIQ_L10C3_FULL_1171
% Mức: VDC
Kéo ngang vật $4\,\text{kg}$ bằng lực 29 $N$ trên sàn có hệ số ma sát 0,2. Lấy g= $10\,\text{m}$ /s^2. Tính gia tốc.
\choice
{$7,25 m/s^2$}
{\True $5,25 m/s^2$}
{$6,3 m/s^2$}
{$4,2 m/s^2$}
\loigiai{
$f=\mu mg=8N$; $a=(F-f)/m=5,25$ m/s^2. Chọn B.
}
\end{ex}

% ===== Câu 38 =====
% ID: AIQ_L10C3_FULL_1172
% Mức: NB
\begin{ex}
% ID: AIQ_L10C3_FULL_1172
% Mức: NB
Kéo ngang vật $5\,\text{kg}$ bằng lực 23 $N$ trên sàn có hệ số ma sát 0,25. Lấy g= $10\,\text{m}$ /s^2. Tính gia tốc.
\choiceTF
{\True Kết quả đúng là $2,1\,\text{m}$ /s^2.}
{Kết quả bằng $4,2\,\text{m}$ /s^2.}
{\True Khi hợp lực bằng 0, vật có gia tốc bằng 0.}
{Có thể bỏ qua đơn vị của đại lượng trong kết quả vật lí.}
\loigiai{
\begin{itemchoice}
\item (Đúng) Kết quả đúng là $2,1\,\text{m}$ /s^2. (Vì): $f=\mu mg=12,5N$; $a=(F-f)/m=2,1$ m/s^2. b) Sai vì giá trị hoặc chiều nêu ra không phù hợp. c) Đúng theo định luật II Newton. d) Sai vì kết quả vật lí phải kèm đơn vị thích hợp
\item (Sai) Kết quả bằng $4,2\,\text{m}$ /s^2.
\item (Đúng) Khi hợp lực bằng 0, vật có gia tốc bằng 0.
\item (Sai) Có thể bỏ qua đơn vị của đại lượng trong kết quả vật lí.
\end{itemchoice}
}
\end{ex}

% ===== Câu 39 =====
% ID: AIQ_L10C3_FULL_1173
% Mức: NB
\begin{ex}
% ID: AIQ_L10C3_FULL_1173
% Mức: NB
Kéo ngang vật $6\,\text{kg}$ bằng lực 25 $N$ trên sàn có hệ số ma sát 0,3. Lấy g= $10\,\text{m}$ /s^2. Tính gia tốc.
\choiceTF
{\True Kết quả đúng là $1,17\,\text{m}$ /s^2.}
{Kết quả bằng $2,34\,\text{m}$ /s^2.}
{\True Khi hợp lực bằng 0, vật có gia tốc bằng 0.}
{Có thể bỏ qua đơn vị của đại lượng trong kết quả vật lí.}
\loigiai{
\begin{itemchoice}
\item (Đúng) Kết quả đúng là $1,17\,\text{m}$ /s^2. (Vì): $f=\mu mg=18N$; $a=(F-f)/m=1,17$ m/s^2. b) Sai vì giá trị hoặc chiều nêu ra không phù hợp. c) Đúng theo định luật II Newton. d) Sai vì kết quả vật lí phải kèm đơn vị thích hợp
\item (Sai) Kết quả bằng $2,34\,\text{m}$ /s^2.
\item (Đúng) Khi hợp lực bằng 0, vật có gia tốc bằng 0.
\item (Sai) Có thể bỏ qua đơn vị của đại lượng trong kết quả vật lí.
\end{itemchoice}
}
\end{ex}

% ===== Câu 40 =====
% ID: AIQ_L10C3_FULL_1174
% Mức: TH
\begin{ex}
% ID: AIQ_L10C3_FULL_1174
% Mức: TH
Kéo ngang vật $7\,\text{kg}$ bằng lực 27 $N$ trên sàn có hệ số ma sát 0,4. Lấy g= $10\,\text{m}$ /s^2. Tính gia tốc.
\choiceTF
{\True Kết quả đúng là - $0,14\,\text{m}$ /s^2.}
{Kết quả bằng - $0,28\,\text{m}$ /s^2.}
{\True Khi hợp lực bằng 0, vật có gia tốc bằng 0.}
{Có thể bỏ qua đơn vị của đại lượng trong kết quả vật lí.}
\loigiai{
\begin{itemchoice}
\item (Đúng) Kết quả đúng là - $0,14\,\text{m}$ /s^2. (Vì): $f=\mu mg=28N$; $a=(F-f)/m=-0,14$ m/s^2. b) Sai vì giá trị hoặc chiều nêu ra không phù hợp. c) Đúng theo định luật II Newton. d) Sai vì kết quả vật lí phải kèm đơn vị thích hợp
\item (Sai) Kết quả bằng - $0,28\,\text{m}$ /s^2.
\item (Đúng) Khi hợp lực bằng 0, vật có gia tốc bằng 0.
\item (Sai) Có thể bỏ qua đơn vị của đại lượng trong kết quả vật lí.
\end{itemchoice}
}
\end{ex}

% ===== Câu 41 =====
% ID: AIQ_L10C3_FULL_1175
% Mức: TH
\begin{ex}
% ID: AIQ_L10C3_FULL_1175
% Mức: TH
Kéo ngang vật $8\,\text{kg}$ bằng lực 29 $N$ trên sàn có hệ số ma sát 0,1. Lấy g= $10\,\text{m}$ /s^2. Tính gia tốc.
\choiceTF
{\True Kết quả đúng là $2,63\,\text{m}$ /s^2.}
{Kết quả bằng $5,26\,\text{m}$ /s^2.}
{\True Khi hợp lực bằng 0, vật có gia tốc bằng 0.}
{Có thể bỏ qua đơn vị của đại lượng trong kết quả vật lí.}
\loigiai{
\begin{itemchoice}
\item (Đúng) Kết quả đúng là $2,63\,\text{m}$ /s^2. (Vì): $f=\mu mg=8N$; $a=(F-f)/m=2,63$ m/s^2. b) Sai vì giá trị hoặc chiều nêu ra không phù hợp. c) Đúng theo định luật II Newton. d) Sai vì kết quả vật lí phải kèm đơn vị thích hợp
\item (Sai) Kết quả bằng $5,26\,\text{m}$ /s^2.
\item (Đúng) Khi hợp lực bằng 0, vật có gia tốc bằng 0.
\item (Sai) Có thể bỏ qua đơn vị của đại lượng trong kết quả vật lí.
\end{itemchoice}
}
\end{ex}

% ===== Câu 42 =====
% ID: AIQ_L10C3_FULL_1176
% Mức: VD
\begin{ex}
% ID: AIQ_L10C3_FULL_1176
% Mức: VD
Kéo ngang vật $2\,\text{kg}$ bằng lực 23 $N$ trên sàn có hệ số ma sát 0,2. Lấy g= $10\,\text{m}$ /s^2. Tính gia tốc.
\choiceTF
{\True Kết quả đúng là $9,5\,\text{m}$ /s^2.}
{Kết quả bằng $19\,\text{m}$ /s^2.}
{\True Khi hợp lực bằng 0, vật có gia tốc bằng 0.}
{Có thể bỏ qua đơn vị của đại lượng trong kết quả vật lí.}
\loigiai{
\begin{itemchoice}
\item (Đúng) Kết quả đúng là $9,5\,\text{m}$ /s^2. (Vì): $f=\mu mg=4N$; $a=(F-f)/m=9,5$ m/s^2. b) Sai vì giá trị hoặc chiều nêu ra không phù hợp. c) Đúng theo định luật II Newton. d) Sai vì kết quả vật lí phải kèm đơn vị thích hợp
\item (Sai) Kết quả bằng $19\,\text{m}$ /s^2.
\item (Đúng) Khi hợp lực bằng 0, vật có gia tốc bằng 0.
\item (Sai) Có thể bỏ qua đơn vị của đại lượng trong kết quả vật lí.
\end{itemchoice}
}
\end{ex}

% ===== Câu 43 =====
% ID: AIQ_L10C3_FULL_1177
% Mức: TH
\begin{ex}
% ID: AIQ_L10C3_FULL_1177
% Mức: TH
Kéo ngang vật $3\,\text{kg}$ bằng lực 25 $N$ trên sàn có hệ số ma sát 0,25. Lấy g= $10\,\text{m}$ /s^2. Tính gia tốc.
\shortans{5,83}
\loigiai{
$f=\mu mg=7,5N$; $a=(F-f)/m=5,83$ m/s^2. Đáp án trả lời ngắn không ghi đơn vị.
}
\end{ex}

% ===== Câu 44 =====
% ID: AIQ_L10C3_FULL_1178
% Mức: TH
\begin{ex}
% ID: AIQ_L10C3_FULL_1178
% Mức: TH
Kéo ngang vật $4\,\text{kg}$ bằng lực 27 $N$ trên sàn có hệ số ma sát 0,3. Lấy g= $10\,\text{m}$ /s^2. Tính gia tốc.
\shortans{3,75}
\loigiai{
$f=\mu mg=12N$; $a=(F-f)/m=3,75$ m/s^2. Đáp án trả lời ngắn không ghi đơn vị.
}
\end{ex}

% ===== Câu 45 =====
% ID: AIQ_L10C3_FULL_1179
% Mức: VD
\begin{ex}
% ID: AIQ_L10C3_FULL_1179
% Mức: VD
Kéo ngang vật $5\,\text{kg}$ bằng lực 29 $N$ trên sàn có hệ số ma sát 0,4. Lấy g= $10\,\text{m}$ /s^2. Tính gia tốc.
\shortans{1,8}
\loigiai{
$f=\mu mg=20N$; $a=(F-f)/m=1,8$ m/s^2. Đáp án trả lời ngắn không ghi đơn vị.
}
\end{ex}

% ===== Câu 46 =====
% ID: AIQ_L10C3_FULL_1180
% Mức: VD
\begin{ex}
% ID: AIQ_L10C3_FULL_1180
% Mức: VD
Kéo ngang vật $6\,\text{kg}$ bằng lực 23 $N$ trên sàn có hệ số ma sát 0,1. Lấy g= $10\,\text{m}$ /s^2. Tính gia tốc.
\shortans{2,83}
\loigiai{
$f=\mu mg=6N$; $a=(F-f)/m=2,83$ m/s^2. Đáp án trả lời ngắn không ghi đơn vị.
}
\end{ex}

% ===== Câu 47 =====
% ID: AIQ_L10C3_FULL_1181
% Mức: VD
\begin{ex}
% ID: AIQ_L10C3_FULL_1181
% Mức: VD
Kéo ngang vật $7\,\text{kg}$ bằng lực 25 $N$ trên sàn có hệ số ma sát 0,2. Lấy g= $10\,\text{m}$ /s^2. Tính gia tốc.
\shortans{1,57}
\loigiai{
$f=\mu mg=14N$; $a=(F-f)/m=1,57$ m/s^2. Đáp án trả lời ngắn không ghi đơn vị.
}
\end{ex}

% ===== Câu 48 =====
% ID: AIQ_L10C3_FULL_1182
% Mức: VDC
\begin{ex}
% ID: AIQ_L10C3_FULL_1182
% Mức: VDC
Kéo ngang vật $8\,\text{kg}$ bằng lực 27 $N$ trên sàn có hệ số ma sát 0,25. Lấy g= $10\,\text{m}$ /s^2. Tính gia tốc.
\shortans{0,88}
\loigiai{
$f=\mu mg=20N$; $a=(F-f)/m=0,88$ m/s^2. Đáp án trả lời ngắn không ghi đơn vị.
}
\end{ex}

% ===== Câu 49 =====
% ID: AIQ_L10C3_FULL_1183
% Mức: TH
\begin{bt}
% ID: AIQ_L10C3_FULL_1183
% Mức: TH
Kéo ngang vật $2\,\text{kg}$ bằng lực 29 $N$ trên sàn có hệ số ma sát 0,3. Lấy g= $10\,\text{m}$ /s^2. Tính gia tốc.
\loigiai{
Phân tích dữ kiện, chọn định luật phù hợp. $f=\mu mg=6N$; $a=(F-f)/m=11,5$ m/s^2.
}
\end{bt}

% ===== Câu 50 =====
% ID: AIQ_L10C3_FULL_1184
% Mức: TH
\begin{bt}
% ID: AIQ_L10C3_FULL_1184
% Mức: TH
Kéo ngang vật $3\,\text{kg}$ bằng lực 23 $N$ trên sàn có hệ số ma sát 0,4. Lấy g= $10\,\text{m}$ /s^2. Tính gia tốc.
\loigiai{
Phân tích dữ kiện, chọn định luật phù hợp. $f=\mu mg=12N$; $a=(F-f)/m=3,67$ m/s^2.
}
\end{bt}

% ===== Câu 51 =====
% ID: AIQ_L10C3_FULL_1185
% Mức: VD
\begin{bt}
% ID: AIQ_L10C3_FULL_1185
% Mức: VD
Kéo ngang vật $4\,\text{kg}$ bằng lực 25 $N$ trên sàn có hệ số ma sát 0,1. Lấy g= $10\,\text{m}$ /s^2. Tính gia tốc.
\loigiai{
Phân tích dữ kiện, chọn định luật phù hợp. $f=\mu mg=4N$; $a=(F-f)/m=5,25$ m/s^2.
}
\end{bt}

% ===== Câu 52 =====
% ID: AIQ_L10C3_FULL_1186
% Mức: VD
\begin{bt}
% ID: AIQ_L10C3_FULL_1186
% Mức: VD
Kéo ngang vật $5\,\text{kg}$ bằng lực 27 $N$ trên sàn có hệ số ma sát 0,2. Lấy g= $10\,\text{m}$ /s^2. Tính gia tốc.
\loigiai{
Phân tích dữ kiện, chọn định luật phù hợp. $f=\mu mg=10N$; $a=(F-f)/m=3,4$ m/s^2.
}
\end{bt}

% ===== Câu 53 =====
% ID: AIQ_L10C3_FULL_1187
% Mức: VD
\begin{bt}
% ID: AIQ_L10C3_FULL_1187
% Mức: VD
Kéo ngang vật $6\,\text{kg}$ bằng lực 29 $N$ trên sàn có hệ số ma sát 0,25. Lấy g= $10\,\text{m}$ /s^2. Tính gia tốc.
\loigiai{
Phân tích dữ kiện, chọn định luật phù hợp. $f=\mu mg=15N$; $a=(F-f)/m=2,33$ m/s^2.
}
\end{bt}

% ===== Câu 54 =====
% ID: AIQ_L10C3_FULL_1188
% Mức: VDC
\begin{bt}
% ID: AIQ_L10C3_FULL_1188
% Mức: VDC
Kéo ngang vật $7\,\text{kg}$ bằng lực 23 $N$ trên sàn có hệ số ma sát 0,3. Lấy g= $10\,\text{m}$ /s^2. Tính gia tốc.
\loigiai{
Phân tích dữ kiện, chọn định luật phù hợp. $f=\mu mg=21N$; $a=(F-f)/m=0,29$ m/s^2.
}
\end{bt}

% ===== Câu 55 =====
% ID: AIQ_L10C3_FULL_1189
% Mức: NB
\dangbt{Chuyển động của vật trên mặt phẳng nghiêng}
\begin{ex}
% ID: AIQ_L10C3_FULL_1189
% Mức: NB
Vật trượt xuống mặt phẳng nghiêng góc 37°, hệ số ma sát 0,4. Lấy g= $10\,\text{m}$ /s^2. Tính gia tốc dọc mặt phẳng nghiêng.
\choice
{$3,36 m/s^2$}
{$2,24 m/s^2$}
{$4,8 m/s^2$}
{\True $2,8 m/s^2$}
\loigiai{
$a=g(\sin\alpha-\mu\cos\alpha)=2,8$ m/s^2. Chọn D.
}
\end{ex}

% ===== Câu 56 =====
% ID: AIQ_L10C3_FULL_1190
% Mức: NB
\begin{ex}
% ID: AIQ_L10C3_FULL_1190
% Mức: NB
Vật trượt xuống mặt phẳng nghiêng góc 45°, hệ số ma sát 0,1. Lấy g= $10\,\text{m}$ /s^2. Tính gia tốc dọc mặt phẳng nghiêng.
\choice
{$5,09 m/s^2$}
{$8,36 m/s^2$}
{\True $6,36 m/s^2$}
{$7,63 m/s^2$}
\loigiai{
$a=g(\sin\alpha-\mu\cos\alpha)=6,36$ m/s^2. Chọn C.
}
\end{ex}

% ===== Câu 57 =====
% ID: AIQ_L10C3_FULL_1191
% Mức: NB
\begin{ex}
% ID: AIQ_L10C3_FULL_1191
% Mức: NB
Vật trượt xuống mặt phẳng nghiêng góc 30°, hệ số ma sát 0,2. Lấy g= $10\,\text{m}$ /s^2. Tính gia tốc dọc mặt phẳng nghiêng.
\choice
{$5,27 m/s^2$}
{\True $3,27 m/s^2$}
{$3,92 m/s^2$}
{$2,62 m/s^2$}
\loigiai{
$a=g(\sin\alpha-\mu\cos\alpha)=3,27$ m/s^2. Chọn B.
}
\end{ex}

% ===== Câu 58 =====
% ID: AIQ_L10C3_FULL_1192
% Mức: TH
\begin{ex}
% ID: AIQ_L10C3_FULL_1192
% Mức: TH
Vật trượt xuống mặt phẳng nghiêng góc 37°, hệ số ma sát 0,25. Lấy g= $10\,\text{m}$ /s^2. Tính gia tốc dọc mặt phẳng nghiêng.
\choice
{\True $4 m/s^2$}
{$4,8 m/s^2$}
{$3,2 m/s^2$}
{$6 m/s^2$}
\loigiai{
$a=g(\sin\alpha-\mu\cos\alpha)=4$ m/s^2. Chọn A.
}
\end{ex}

% ===== Câu 59 =====
% ID: AIQ_L10C3_FULL_1193
% Mức: TH
\begin{ex}
% ID: AIQ_L10C3_FULL_1193
% Mức: TH
Vật trượt xuống mặt phẳng nghiêng góc 45°, hệ số ma sát 0,3. Lấy g= $10\,\text{m}$ /s^2. Tính gia tốc dọc mặt phẳng nghiêng.
\choice
{$5,94 m/s^2$}
{$3,96 m/s^2$}
{$6,95 m/s^2$}
{\True $4,95 m/s^2$}
\loigiai{
$a=g(\sin\alpha-\mu\cos\alpha)=4,95$ m/s^2. Chọn D.
}
\end{ex}

% ===== Câu 60 =====
% ID: AIQ_L10C3_FULL_1194
% Mức: TH
\begin{ex}
% ID: AIQ_L10C3_FULL_1194
% Mức: TH
Vật trượt xuống mặt phẳng nghiêng góc 30°, hệ số ma sát 0,4. Lấy g= $10\,\text{m}$ /s^2. Tính gia tốc dọc mặt phẳng nghiêng.
\choice
{$1,23 m/s^2$}
{$3,54 m/s^2$}
{\True $1,54 m/s^2$}
{$1,85 m/s^2$}
\loigiai{
$a=g(\sin\alpha-\mu\cos\alpha)=1,54$ m/s^2. Chọn C.
}
\end{ex}

% ===== Câu 61 =====
% ID: AIQ_L10C3_FULL_1195
% Mức: TH
\begin{ex}
% ID: AIQ_L10C3_FULL_1195
% Mức: TH
Vật trượt xuống mặt phẳng nghiêng góc 37°, hệ số ma sát 0,1. Lấy g= $10\,\text{m}$ /s^2. Tính gia tốc dọc mặt phẳng nghiêng.
\choice
{$7,2 m/s^2$}
{\True $5,2 m/s^2$}
{$6,24 m/s^2$}
{$4,16 m/s^2$}
\loigiai{
$a=g(\sin\alpha-\mu\cos\alpha)=5,2$ m/s^2. Chọn B.
}
\end{ex}

% ===== Câu 62 =====
% ID: AIQ_L10C3_FULL_1196
% Mức: VD
\begin{ex}
% ID: AIQ_L10C3_FULL_1196
% Mức: VD
Vật trượt xuống mặt phẳng nghiêng góc 45°, hệ số ma sát 0,2. Lấy g= $10\,\text{m}$ /s^2. Tính gia tốc dọc mặt phẳng nghiêng.
\choice
{\True $5,66 m/s^2$}
{$6,79 m/s^2$}
{$4,53 m/s^2$}
{$7,66 m/s^2$}
\loigiai{
$a=g(\sin\alpha-\mu\cos\alpha)=5,66$ m/s^2. Chọn A.
}
\end{ex}

% ===== Câu 63 =====
% ID: AIQ_L10C3_FULL_1197
% Mức: VD
\begin{ex}
% ID: AIQ_L10C3_FULL_1197
% Mức: VD
Vật trượt xuống mặt phẳng nghiêng góc 30°, hệ số ma sát 0,25. Lấy g= $10\,\text{m}$ /s^2. Tính gia tốc dọc mặt phẳng nghiêng.
\choice
{$3,41 m/s^2$}
{$2,27 m/s^2$}
{$4,84 m/s^2$}
{\True $2,84 m/s^2$}
\loigiai{
$a=g(\sin\alpha-\mu\cos\alpha)=2,84$ m/s^2. Chọn D.
}
\end{ex}

% ===== Câu 64 =====
% ID: AIQ_L10C3_FULL_1198
% Mức: VDC
\begin{ex}
% ID: AIQ_L10C3_FULL_1198
% Mức: VDC
Vật trượt xuống mặt phẳng nghiêng góc 37°, hệ số ma sát 0,3. Lấy g= $10\,\text{m}$ /s^2. Tính gia tốc dọc mặt phẳng nghiêng.
\choice
{$2,88 m/s^2$}
{$5,6 m/s^2$}
{\True $3,6 m/s^2$}
{$4,32 m/s^2$}
\loigiai{
$a=g(\sin\alpha-\mu\cos\alpha)=3,6$ m/s^2. Chọn C.
}
\end{ex}

% ===== Câu 65 =====
% ID: AIQ_L10C3_FULL_1199
% Mức: NB
\begin{ex}
% ID: AIQ_L10C3_FULL_1199
% Mức: NB
Vật trượt xuống mặt phẳng nghiêng góc 45°, hệ số ma sát 0,4. Lấy g= $10\,\text{m}$ /s^2. Tính gia tốc dọc mặt phẳng nghiêng.
\choiceTF
{\True Kết quả đúng là $4,24\,\text{m}$ /s^2.}
{Kết quả bằng $8,48\,\text{m}$ /s^2.}
{\True Khi hợp lực bằng 0, vật có gia tốc bằng 0.}
{Có thể bỏ qua đơn vị của đại lượng trong kết quả vật lí.}
\loigiai{
\begin{itemchoice}
\item (Đúng) Kết quả đúng là $4,24\,\text{m}$ /s^2. (Vì): $a=g(\sin\alpha-\mu\cos\alpha)=4,24$ m/s^2. b) Sai vì giá trị hoặc chiều nêu ra không phù hợp. c) Đúng theo định luật II Newton. d) Sai vì kết quả vật lí phải kèm đơn vị thích hợp
\item (Sai) Kết quả bằng $8,48\,\text{m}$ /s^2.
\item (Đúng) Khi hợp lực bằng 0, vật có gia tốc bằng 0.
\item (Sai) Có thể bỏ qua đơn vị của đại lượng trong kết quả vật lí.
\end{itemchoice}
}
\end{ex}

% ===== Câu 66 =====
% ID: AIQ_L10C3_FULL_1200
% Mức: NB
\begin{ex}
% ID: AIQ_L10C3_FULL_1200
% Mức: NB
Vật trượt xuống mặt phẳng nghiêng góc 30°, hệ số ma sát 0,1. Lấy g= $10\,\text{m}$ /s^2. Tính gia tốc dọc mặt phẳng nghiêng.
\choiceTF
{\True Kết quả đúng là $4,13\,\text{m}$ /s^2.}
{Kết quả bằng $8,26\,\text{m}$ /s^2.}
{\True Khi hợp lực bằng 0, vật có gia tốc bằng 0.}
{Có thể bỏ qua đơn vị của đại lượng trong kết quả vật lí.}
\loigiai{
\begin{itemchoice}
\item (Đúng) Kết quả đúng là $4,13\,\text{m}$ /s^2. (Vì): $a=g(\sin\alpha-\mu\cos\alpha)=4,13$ m/s^2. b) Sai vì giá trị hoặc chiều nêu ra không phù hợp. c) Đúng theo định luật II Newton. d) Sai vì kết quả vật lí phải kèm đơn vị thích hợp
\item (Sai) Kết quả bằng $8,26\,\text{m}$ /s^2.
\item (Đúng) Khi hợp lực bằng 0, vật có gia tốc bằng 0.
\item (Sai) Có thể bỏ qua đơn vị của đại lượng trong kết quả vật lí.
\end{itemchoice}
}
\end{ex}

% ===== Câu 67 =====
% ID: AIQ_L10C3_FULL_1201
% Mức: TH
\begin{ex}
% ID: AIQ_L10C3_FULL_1201
% Mức: TH
Vật trượt xuống mặt phẳng nghiêng góc 37°, hệ số ma sát 0,2. Lấy g= $10\,\text{m}$ /s^2. Tính gia tốc dọc mặt phẳng nghiêng.
\choiceTF
{\True Kết quả đúng là $4,4\,\text{m}$ /s^2.}
{Kết quả bằng $8,8\,\text{m}$ /s^2.}
{\True Khi hợp lực bằng 0, vật có gia tốc bằng 0.}
{Có thể bỏ qua đơn vị của đại lượng trong kết quả vật lí.}
\loigiai{
\begin{itemchoice}
\item (Đúng) Kết quả đúng là $4,4\,\text{m}$ /s^2. (Vì): $a=g(\sin\alpha-\mu\cos\alpha)=4,4$ m/s^2. b) Sai vì giá trị hoặc chiều nêu ra không phù hợp. c) Đúng theo định luật II Newton. d) Sai vì kết quả vật lí phải kèm đơn vị thích hợp
\item (Sai) Kết quả bằng $8,8\,\text{m}$ /s^2.
\item (Đúng) Khi hợp lực bằng 0, vật có gia tốc bằng 0.
\item (Sai) Có thể bỏ qua đơn vị của đại lượng trong kết quả vật lí.
\end{itemchoice}
}
\end{ex}

% ===== Câu 68 =====
% ID: AIQ_L10C3_FULL_1202
% Mức: TH
\begin{ex}
% ID: AIQ_L10C3_FULL_1202
% Mức: TH
Vật trượt xuống mặt phẳng nghiêng góc 45°, hệ số ma sát 0,25. Lấy g= $10\,\text{m}$ /s^2. Tính gia tốc dọc mặt phẳng nghiêng.
\choiceTF
{\True Kết quả đúng là $5,3\,\text{m}$ /s^2.}
{Kết quả bằng $10,6\,\text{m}$ /s^2.}
{\True Khi hợp lực bằng 0, vật có gia tốc bằng 0.}
{Có thể bỏ qua đơn vị của đại lượng trong kết quả vật lí.}
\loigiai{
\begin{itemchoice}
\item (Đúng) Kết quả đúng là $5,3\,\text{m}$ /s^2. (Vì): $a=g(\sin\alpha-\mu\cos\alpha)=5,3$ m/s^2. b) Sai vì giá trị hoặc chiều nêu ra không phù hợp. c) Đúng theo định luật II Newton. d) Sai vì kết quả vật lí phải kèm đơn vị thích hợp
\item (Sai) Kết quả bằng $10,6\,\text{m}$ /s^2.
\item (Đúng) Khi hợp lực bằng 0, vật có gia tốc bằng 0.
\item (Sai) Có thể bỏ qua đơn vị của đại lượng trong kết quả vật lí.
\end{itemchoice}
}
\end{ex}

% ===== Câu 69 =====
% ID: AIQ_L10C3_FULL_1203
% Mức: VD
\begin{ex}
% ID: AIQ_L10C3_FULL_1203
% Mức: VD
Vật trượt xuống mặt phẳng nghiêng góc 30°, hệ số ma sát 0,3. Lấy g= $10\,\text{m}$ /s^2. Tính gia tốc dọc mặt phẳng nghiêng.
\choiceTF
{\True Kết quả đúng là $2,4\,\text{m}$ /s^2.}
{Kết quả bằng $4,8\,\text{m}$ /s^2.}
{\True Khi hợp lực bằng 0, vật có gia tốc bằng 0.}
{Có thể bỏ qua đơn vị của đại lượng trong kết quả vật lí.}
\loigiai{
\begin{itemchoice}
\item (Đúng) Kết quả đúng là $2,4\,\text{m}$ /s^2. (Vì): $a=g(\sin\alpha-\mu\cos\alpha)=2,4$ m/s^2. b) Sai vì giá trị hoặc chiều nêu ra không phù hợp. c) Đúng theo định luật II Newton. d) Sai vì kết quả vật lí phải kèm đơn vị thích hợp
\item (Sai) Kết quả bằng $4,8\,\text{m}$ /s^2.
\item (Đúng) Khi hợp lực bằng 0, vật có gia tốc bằng 0.
\item (Sai) Có thể bỏ qua đơn vị của đại lượng trong kết quả vật lí.
\end{itemchoice}
}
\end{ex}

% ===== Câu 70 =====
% ID: AIQ_L10C3_FULL_1204
% Mức: TH
\begin{ex}
% ID: AIQ_L10C3_FULL_1204
% Mức: TH
Vật trượt xuống mặt phẳng nghiêng góc 37°, hệ số ma sát 0,4. Lấy g= $10\,\text{m}$ /s^2. Tính gia tốc dọc mặt phẳng nghiêng.
\shortans{2,8}
\loigiai{
$a=g(\sin\alpha-\mu\cos\alpha)=2,8$ m/s^2. Đáp án trả lời ngắn không ghi đơn vị.
}
\end{ex}

% ===== Câu 71 =====
% ID: AIQ_L10C3_FULL_1205
% Mức: TH
\begin{ex}
% ID: AIQ_L10C3_FULL_1205
% Mức: TH
Vật trượt xuống mặt phẳng nghiêng góc 45°, hệ số ma sát 0,1. Lấy g= $10\,\text{m}$ /s^2. Tính gia tốc dọc mặt phẳng nghiêng.
\shortans{6,36}
\loigiai{
$a=g(\sin\alpha-\mu\cos\alpha)=6,36$ m/s^2. Đáp án trả lời ngắn không ghi đơn vị.
}
\end{ex}

% ===== Câu 72 =====
% ID: AIQ_L10C3_FULL_1206
% Mức: VD
\begin{ex}
% ID: AIQ_L10C3_FULL_1206
% Mức: VD
Vật trượt xuống mặt phẳng nghiêng góc 30°, hệ số ma sát 0,2. Lấy g= $10\,\text{m}$ /s^2. Tính gia tốc dọc mặt phẳng nghiêng.
\shortans{3,27}
\loigiai{
$a=g(\sin\alpha-\mu\cos\alpha)=3,27$ m/s^2. Đáp án trả lời ngắn không ghi đơn vị.
}
\end{ex}

% ===== Câu 73 =====
% ID: AIQ_L10C3_FULL_1207
% Mức: VD
\begin{ex}
% ID: AIQ_L10C3_FULL_1207
% Mức: VD
Vật trượt xuống mặt phẳng nghiêng góc 37°, hệ số ma sát 0,25. Lấy g= $10\,\text{m}$ /s^2. Tính gia tốc dọc mặt phẳng nghiêng.
\shortans{4}
\loigiai{
$a=g(\sin\alpha-\mu\cos\alpha)=4$ m/s^2. Đáp án trả lời ngắn không ghi đơn vị.
}
\end{ex}

% ===== Câu 74 =====
% ID: AIQ_L10C3_FULL_1208
% Mức: VD
\begin{ex}
% ID: AIQ_L10C3_FULL_1208
% Mức: VD
Vật trượt xuống mặt phẳng nghiêng góc 45°, hệ số ma sát 0,3. Lấy g= $10\,\text{m}$ /s^2. Tính gia tốc dọc mặt phẳng nghiêng.
\shortans{4,95}
\loigiai{
$a=g(\sin\alpha-\mu\cos\alpha)=4,95$ m/s^2. Đáp án trả lời ngắn không ghi đơn vị.
}
\end{ex}

% ===== Câu 75 =====
% ID: AIQ_L10C3_FULL_1209
% Mức: VDC
\begin{ex}
% ID: AIQ_L10C3_FULL_1209
% Mức: VDC
Vật trượt xuống mặt phẳng nghiêng góc 30°, hệ số ma sát 0,4. Lấy g= $10\,\text{m}$ /s^2. Tính gia tốc dọc mặt phẳng nghiêng.
\shortans{1,54}
\loigiai{
$a=g(\sin\alpha-\mu\cos\alpha)=1,54$ m/s^2. Đáp án trả lời ngắn không ghi đơn vị.
}
\end{ex}

% ===== Câu 76 =====
% ID: AIQ_L10C3_FULL_1210
% Mức: TH
\begin{bt}
% ID: AIQ_L10C3_FULL_1210
% Mức: TH
Vật trượt xuống mặt phẳng nghiêng góc 37°, hệ số ma sát 0,1. Lấy g= $10\,\text{m}$ /s^2. Tính gia tốc dọc mặt phẳng nghiêng.
\loigiai{
Phân tích dữ kiện, chọn định luật phù hợp. $a=g(\sin\alpha-\mu\cos\alpha)=5,2$ m/s^2.
}
\end{bt}

% ===== Câu 77 =====
% ID: AIQ_L10C3_FULL_1211
% Mức: TH
\begin{bt}
% ID: AIQ_L10C3_FULL_1211
% Mức: TH
Vật trượt xuống mặt phẳng nghiêng góc 45°, hệ số ma sát 0,2. Lấy g= $10\,\text{m}$ /s^2. Tính gia tốc dọc mặt phẳng nghiêng.
\loigiai{
Phân tích dữ kiện, chọn định luật phù hợp. $a=g(\sin\alpha-\mu\cos\alpha)=5,66$ m/s^2.
}
\end{bt}

% ===== Câu 78 =====
% ID: AIQ_L10C3_FULL_1212
% Mức: VD
\begin{bt}
% ID: AIQ_L10C3_FULL_1212
% Mức: VD
Vật trượt xuống mặt phẳng nghiêng góc 30°, hệ số ma sát 0,25. Lấy g= $10\,\text{m}$ /s^2. Tính gia tốc dọc mặt phẳng nghiêng.
\loigiai{
Phân tích dữ kiện, chọn định luật phù hợp. $a=g(\sin\alpha-\mu\cos\alpha)=2,84$ m/s^2.
}
\end{bt}

% ===== Câu 79 =====
% ID: AIQ_L10C3_FULL_1213
% Mức: VD
\begin{bt}
% ID: AIQ_L10C3_FULL_1213
% Mức: VD
Vật trượt xuống mặt phẳng nghiêng góc 37°, hệ số ma sát 0,3. Lấy g= $10\,\text{m}$ /s^2. Tính gia tốc dọc mặt phẳng nghiêng.
\loigiai{
Phân tích dữ kiện, chọn định luật phù hợp. $a=g(\sin\alpha-\mu\cos\alpha)=3,6$ m/s^2.
}
\end{bt}

% ===== Câu 80 =====
% ID: AIQ_L10C3_FULL_1214
% Mức: VD
\begin{bt}
% ID: AIQ_L10C3_FULL_1214
% Mức: VD
Vật trượt xuống mặt phẳng nghiêng góc 45°, hệ số ma sát 0,4. Lấy g= $10\,\text{m}$ /s^2. Tính gia tốc dọc mặt phẳng nghiêng.
\loigiai{
Phân tích dữ kiện, chọn định luật phù hợp. $a=g(\sin\alpha-\mu\cos\alpha)=4,24$ m/s^2.
}
\end{bt}

% ===== Câu 81 =====
% ID: AIQ_L10C3_FULL_1215
% Mức: VDC
\begin{bt}
% ID: AIQ_L10C3_FULL_1215
% Mức: VDC
Vật trượt xuống mặt phẳng nghiêng góc 30°, hệ số ma sát 0,1. Lấy g= $10\,\text{m}$ /s^2. Tính gia tốc dọc mặt phẳng nghiêng.
\loigiai{
Phân tích dữ kiện, chọn định luật phù hợp. $a=g(\sin\alpha-\mu\cos\alpha)=4,13$ m/s^2.
}
\end{bt}

% ===== Câu 82 =====
% ID: AIQ_L10C3_FULL_1216
% Mức: NB
\dangbt{Hệ hai vật nối với nhau bằng dây không dãn}
\begin{ex}
% ID: AIQ_L10C3_FULL_1216
% Mức: NB
Vật $7\,\text{kg}$ trên bàn ngang có ma sát $\mu=0,2$ nối qua ròng rọc với vật treo $9\,\text{kg}$. Bỏ qua khối lượng dây, ròng rọc. Lấy g= $10\,\text{m}$ /s^2. Tính gia tốc hệ.
\choice
{\True $4,75 m/s^2$}
{$5,7 m/s^2$}
{$3,8 m/s^2$}
{$6,75 m/s^2$}
\loigiai{
Với toàn hệ: $a=(m_2g-\mu m_1g)/(m_1+m_2)=4,75$ m/s^2. Chọn A.
}
\end{ex}

% ===== Câu 83 =====
% ID: AIQ_L10C3_FULL_1217
% Mức: NB
\begin{ex}
% ID: AIQ_L10C3_FULL_1217
% Mức: NB
Vật $8\,\text{kg}$ trên bàn ngang có ma sát $\mu=0,25$ nối qua ròng rọc với vật treo $10\,\text{kg}$. Bỏ qua khối lượng dây, ròng rọc. Lấy g= $10\,\text{m}$ /s^2. Tính gia tốc hệ.
\choice
{$5,33 m/s^2$}
{$3,55 m/s^2$}
{$6,44 m/s^2$}
{\True $4,44 m/s^2$}
\loigiai{
Với toàn hệ: $a=(m_2g-\mu m_1g)/(m_1+m_2)=4,44$ m/s^2. Chọn D.
}
\end{ex}

% ===== Câu 84 =====
% ID: AIQ_L10C3_FULL_1218
% Mức: NB
\begin{ex}
% ID: AIQ_L10C3_FULL_1218
% Mức: NB
Vật $2\,\text{kg}$ trên bàn ngang có ma sát $\mu=0,3$ nối qua ròng rọc với vật treo $4\,\text{kg}$. Bỏ qua khối lượng dây, ròng rọc. Lấy g= $10\,\text{m}$ /s^2. Tính gia tốc hệ.
\choice
{$4,54 m/s^2$}
{$7,67 m/s^2$}
{\True $5,67 m/s^2$}
{$6,8 m/s^2$}
\loigiai{
Với toàn hệ: $a=(m_2g-\mu m_1g)/(m_1+m_2)=5,67$ m/s^2. Chọn C.
}
\end{ex}

% ===== Câu 85 =====
% ID: AIQ_L10C3_FULL_1219
% Mức: TH
\begin{ex}
% ID: AIQ_L10C3_FULL_1219
% Mức: TH
Vật $3\,\text{kg}$ trên bàn ngang có ma sát $\mu=0,4$ nối qua ròng rọc với vật treo $5\,\text{kg}$. Bỏ qua khối lượng dây, ròng rọc. Lấy g= $10\,\text{m}$ /s^2. Tính gia tốc hệ.
\choice
{$6,75 m/s^2$}
{\True $4,75 m/s^2$}
{$5,7 m/s^2$}
{$3,8 m/s^2$}
\loigiai{
Với toàn hệ: $a=(m_2g-\mu m_1g)/(m_1+m_2)=4,75$ m/s^2. Chọn B.
}
\end{ex}

% ===== Câu 86 =====
% ID: AIQ_L10C3_FULL_1220
% Mức: TH
\begin{ex}
% ID: AIQ_L10C3_FULL_1220
% Mức: TH
Vật $4\,\text{kg}$ trên bàn ngang có ma sát $\mu=0,1$ nối qua ròng rọc với vật treo $6\,\text{kg}$. Bỏ qua khối lượng dây, ròng rọc. Lấy g= $10\,\text{m}$ /s^2. Tính gia tốc hệ.
\choice
{\True $5,6 m/s^2$}
{$6,72 m/s^2$}
{$4,48 m/s^2$}
{$7,6 m/s^2$}
\loigiai{
Với toàn hệ: $a=(m_2g-\mu m_1g)/(m_1+m_2)=5,6$ m/s^2. Chọn A.
}
\end{ex}

% ===== Câu 87 =====
% ID: AIQ_L10C3_FULL_1221
% Mức: TH
\begin{ex}
% ID: AIQ_L10C3_FULL_1221
% Mức: TH
Vật $5\,\text{kg}$ trên bàn ngang có ma sát $\mu=0,2$ nối qua ròng rọc với vật treo $7\,\text{kg}$. Bỏ qua khối lượng dây, ròng rọc. Lấy g= $10\,\text{m}$ /s^2. Tính gia tốc hệ.
\choice
{$6 m/s^2$}
{$4 m/s^2$}
{$7 m/s^2$}
{\True $5 m/s^2$}
\loigiai{
Với toàn hệ: $a=(m_2g-\mu m_1g)/(m_1+m_2)=5$ m/s^2. Chọn D.
}
\end{ex}

% ===== Câu 88 =====
% ID: AIQ_L10C3_FULL_1222
% Mức: TH
\begin{ex}
% ID: AIQ_L10C3_FULL_1222
% Mức: TH
Vật $6\,\text{kg}$ trên bàn ngang có ma sát $\mu=0,25$ nối qua ròng rọc với vật treo $8\,\text{kg}$. Bỏ qua khối lượng dây, ròng rọc. Lấy g= $10\,\text{m}$ /s^2. Tính gia tốc hệ.
\choice
{$3,71 m/s^2$}
{$6,64 m/s^2$}
{\True $4,64 m/s^2$}
{$5,57 m/s^2$}
\loigiai{
Với toàn hệ: $a=(m_2g-\mu m_1g)/(m_1+m_2)=4,64$ m/s^2. Chọn C.
}
\end{ex}

% ===== Câu 89 =====
% ID: AIQ_L10C3_FULL_1223
% Mức: VD
\begin{ex}
% ID: AIQ_L10C3_FULL_1223
% Mức: VD
Vật $7\,\text{kg}$ trên bàn ngang có ma sát $\mu=0,3$ nối qua ròng rọc với vật treo $9\,\text{kg}$. Bỏ qua khối lượng dây, ròng rọc. Lấy g= $10\,\text{m}$ /s^2. Tính gia tốc hệ.
\choice
{$6,31 m/s^2$}
{\True $4,31 m/s^2$}
{$5,17 m/s^2$}
{$3,45 m/s^2$}
\loigiai{
Với toàn hệ: $a=(m_2g-\mu m_1g)/(m_1+m_2)=4,31$ m/s^2. Chọn B.
}
\end{ex}

% ===== Câu 90 =====
% ID: AIQ_L10C3_FULL_1224
% Mức: VD
\begin{ex}
% ID: AIQ_L10C3_FULL_1224
% Mức: VD
Vật $8\,\text{kg}$ trên bàn ngang có ma sát $\mu=0,4$ nối qua ròng rọc với vật treo $10\,\text{kg}$. Bỏ qua khối lượng dây, ròng rọc. Lấy g= $10\,\text{m}$ /s^2. Tính gia tốc hệ.
\choice
{\True $3,78 m/s^2$}
{$4,54 m/s^2$}
{$3,02 m/s^2$}
{$5,78 m/s^2$}
\loigiai{
Với toàn hệ: $a=(m_2g-\mu m_1g)/(m_1+m_2)=3,78$ m/s^2. Chọn A.
}
\end{ex}

% ===== Câu 91 =====
% ID: AIQ_L10C3_FULL_1225
% Mức: VDC
\begin{ex}
% ID: AIQ_L10C3_FULL_1225
% Mức: VDC
Vật $2\,\text{kg}$ trên bàn ngang có ma sát $\mu=0,1$ nối qua ròng rọc với vật treo $4\,\text{kg}$. Bỏ qua khối lượng dây, ròng rọc. Lấy g= $10\,\text{m}$ /s^2. Tính gia tốc hệ.
\choice
{$7,6 m/s^2$}
{$5,06 m/s^2$}
{$8,33 m/s^2$}
{\True $6,33 m/s^2$}
\loigiai{
Với toàn hệ: $a=(m_2g-\mu m_1g)/(m_1+m_2)=6,33$ m/s^2. Chọn D.
}
\end{ex}

% ===== Câu 92 =====
% ID: AIQ_L10C3_FULL_1226
% Mức: NB
\begin{ex}
% ID: AIQ_L10C3_FULL_1226
% Mức: NB
Vật $3\,\text{kg}$ trên bàn ngang có ma sát $\mu=0,2$ nối qua ròng rọc với vật treo $5\,\text{kg}$. Bỏ qua khối lượng dây, ròng rọc. Lấy g= $10\,\text{m}$ /s^2. Tính gia tốc hệ.
\choiceTF
{\True Kết quả đúng là $5,5\,\text{m}$ /s^2.}
{Kết quả bằng $11\,\text{m}$ /s^2.}
{\True Khi hợp lực bằng 0, vật có gia tốc bằng 0.}
{Có thể bỏ qua đơn vị của đại lượng trong kết quả vật lí.}
\loigiai{
\begin{itemchoice}
\item (Đúng) Kết quả đúng là $5,5\,\text{m}$ /s^2. (Vì): Với toàn hệ: $a=(m_2g-\mu m_1g)/(m_1+m_2)=5,5$ m/s^2. b) Sai vì giá trị hoặc chiều nêu ra không phù hợp. c) Đúng theo định luật II Newton. d) Sai vì kết quả vật lí phải kèm đơn vị thích hợp
\item (Sai) Kết quả bằng $11\,\text{m}$ /s^2.
\item (Đúng) Khi hợp lực bằng 0, vật có gia tốc bằng 0.
\item (Sai) Có thể bỏ qua đơn vị của đại lượng trong kết quả vật lí.
\end{itemchoice}
}
\end{ex}

% ===== Câu 93 =====
% ID: AIQ_L10C3_FULL_1227
% Mức: NB
\begin{ex}
% ID: AIQ_L10C3_FULL_1227
% Mức: NB
Vật $4\,\text{kg}$ trên bàn ngang có ma sát $\mu=0,25$ nối qua ròng rọc với vật treo $6\,\text{kg}$. Bỏ qua khối lượng dây, ròng rọc. Lấy g= $10\,\text{m}$ /s^2. Tính gia tốc hệ.
\choiceTF
{\True Kết quả đúng là $5\,\text{m}$ /s^2.}
{Kết quả bằng $10\,\text{m}$ /s^2.}
{\True Khi hợp lực bằng 0, vật có gia tốc bằng 0.}
{Có thể bỏ qua đơn vị của đại lượng trong kết quả vật lí.}
\loigiai{
\begin{itemchoice}
\item (Đúng) Kết quả đúng là $5\,\text{m}$ /s^2. (Vì): Với toàn hệ: $a=(m_2g-\mu m_1g)/(m_1+m_2)=5$ m/s^2. b) Sai vì giá trị hoặc chiều nêu ra không phù hợp. c) Đúng theo định luật II Newton. d) Sai vì kết quả vật lí phải kèm đơn vị thích hợp
\item (Sai) Kết quả bằng $10\,\text{m}$ /s^2.
\item (Đúng) Khi hợp lực bằng 0, vật có gia tốc bằng 0.
\item (Sai) Có thể bỏ qua đơn vị của đại lượng trong kết quả vật lí.
\end{itemchoice}
}
\end{ex}

% ===== Câu 94 =====
% ID: AIQ_L10C3_FULL_1228
% Mức: TH
\begin{ex}
% ID: AIQ_L10C3_FULL_1228
% Mức: TH
Vật $5\,\text{kg}$ trên bàn ngang có ma sát $\mu=0,3$ nối qua ròng rọc với vật treo $7\,\text{kg}$. Bỏ qua khối lượng dây, ròng rọc. Lấy g= $10\,\text{m}$ /s^2. Tính gia tốc hệ.
\choiceTF
{\True Kết quả đúng là $4,58\,\text{m}$ /s^2.}
{Kết quả bằng $9,16\,\text{m}$ /s^2.}
{\True Khi hợp lực bằng 0, vật có gia tốc bằng 0.}
{Có thể bỏ qua đơn vị của đại lượng trong kết quả vật lí.}
\loigiai{
\begin{itemchoice}
\item (Đúng) Kết quả đúng là $4,58\,\text{m}$ /s^2. (Vì): Với toàn hệ: $a=(m_2g-\mu m_1g)/(m_1+m_2)=4,58$ m/s^2. b) Sai vì giá trị hoặc chiều nêu ra không phù hợp. c) Đúng theo định luật II Newton. d) Sai vì kết quả vật lí phải kèm đơn vị thích hợp
\item (Sai) Kết quả bằng $9,16\,\text{m}$ /s^2.
\item (Đúng) Khi hợp lực bằng 0, vật có gia tốc bằng 0.
\item (Sai) Có thể bỏ qua đơn vị của đại lượng trong kết quả vật lí.
\end{itemchoice}
}
\end{ex}

% ===== Câu 95 =====
% ID: AIQ_L10C3_FULL_1229
% Mức: TH
\begin{ex}
% ID: AIQ_L10C3_FULL_1229
% Mức: TH
Vật $6\,\text{kg}$ trên bàn ngang có ma sát $\mu=0,4$ nối qua ròng rọc với vật treo $8\,\text{kg}$. Bỏ qua khối lượng dây, ròng rọc. Lấy g= $10\,\text{m}$ /s^2. Tính gia tốc hệ.
\choiceTF
{\True Kết quả đúng là $4\,\text{m}$ /s^2.}
{Kết quả bằng $8\,\text{m}$ /s^2.}
{\True Khi hợp lực bằng 0, vật có gia tốc bằng 0.}
{Có thể bỏ qua đơn vị của đại lượng trong kết quả vật lí.}
\loigiai{
\begin{itemchoice}
\item (Đúng) Kết quả đúng là $4\,\text{m}$ /s^2. (Vì): Với toàn hệ: $a=(m_2g-\mu m_1g)/(m_1+m_2)=4$ m/s^2. b) Sai vì giá trị hoặc chiều nêu ra không phù hợp. c) Đúng theo định luật II Newton. d) Sai vì kết quả vật lí phải kèm đơn vị thích hợp
\item (Sai) Kết quả bằng $8\,\text{m}$ /s^2.
\item (Đúng) Khi hợp lực bằng 0, vật có gia tốc bằng 0.
\item (Sai) Có thể bỏ qua đơn vị của đại lượng trong kết quả vật lí.
\end{itemchoice}
}
\end{ex}

% ===== Câu 96 =====
% ID: AIQ_L10C3_FULL_1230
% Mức: VD
\begin{ex}
% ID: AIQ_L10C3_FULL_1230
% Mức: VD
Vật $7\,\text{kg}$ trên bàn ngang có ma sát $\mu=0,1$ nối qua ròng rọc với vật treo $9\,\text{kg}$. Bỏ qua khối lượng dây, ròng rọc. Lấy g= $10\,\text{m}$ /s^2. Tính gia tốc hệ.
\choiceTF
{\True Kết quả đúng là $5,19\,\text{m}$ /s^2.}
{Kết quả bằng $10,38\,\text{m}$ /s^2.}
{\True Khi hợp lực bằng 0, vật có gia tốc bằng 0.}
{Có thể bỏ qua đơn vị của đại lượng trong kết quả vật lí.}
\loigiai{
\begin{itemchoice}
\item (Đúng) Kết quả đúng là $5,19\,\text{m}$ /s^2. (Vì): Với toàn hệ: $a=(m_2g-\mu m_1g)/(m_1+m_2)=5,19$ m/s^2. b) Sai vì giá trị hoặc chiều nêu ra không phù hợp. c) Đúng theo định luật II Newton. d) Sai vì kết quả vật lí phải kèm đơn vị thích hợp
\item (Sai) Kết quả bằng $10,38\,\text{m}$ /s^2.
\item (Đúng) Khi hợp lực bằng 0, vật có gia tốc bằng 0.
\item (Sai) Có thể bỏ qua đơn vị của đại lượng trong kết quả vật lí.
\end{itemchoice}
}
\end{ex}

% ===== Câu 97 =====
% ID: AIQ_L10C3_FULL_1231
% Mức: TH
\begin{ex}
% ID: AIQ_L10C3_FULL_1231
% Mức: TH
Vật $8\,\text{kg}$ trên bàn ngang có ma sát $\mu=0,2$ nối qua ròng rọc với vật treo $10\,\text{kg}$. Bỏ qua khối lượng dây, ròng rọc. Lấy g= $10\,\text{m}$ /s^2. Tính gia tốc hệ.
\shortans{4,67}
\loigiai{
Với toàn hệ: $a=(m_2g-\mu m_1g)/(m_1+m_2)=4,67$ m/s^2. Đáp án trả lời ngắn không ghi đơn vị.
}
\end{ex}

% ===== Câu 98 =====
% ID: AIQ_L10C3_FULL_1232
% Mức: TH
\begin{ex}
% ID: AIQ_L10C3_FULL_1232
% Mức: TH
Vật $2\,\text{kg}$ trên bàn ngang có ma sát $\mu=0,25$ nối qua ròng rọc với vật treo $4\,\text{kg}$. Bỏ qua khối lượng dây, ròng rọc. Lấy g= $10\,\text{m}$ /s^2. Tính gia tốc hệ.
\shortans{5,83}
\loigiai{
Với toàn hệ: $a=(m_2g-\mu m_1g)/(m_1+m_2)=5,83$ m/s^2. Đáp án trả lời ngắn không ghi đơn vị.
}
\end{ex}

% ===== Câu 99 =====
% ID: AIQ_L10C3_FULL_1233
% Mức: VD
\begin{ex}
% ID: AIQ_L10C3_FULL_1233
% Mức: VD
Vật $3\,\text{kg}$ trên bàn ngang có ma sát $\mu=0,3$ nối qua ròng rọc với vật treo $5\,\text{kg}$. Bỏ qua khối lượng dây, ròng rọc. Lấy g= $10\,\text{m}$ /s^2. Tính gia tốc hệ.
\shortans{5,13}
\loigiai{
Với toàn hệ: $a=(m_2g-\mu m_1g)/(m_1+m_2)=5,13$ m/s^2. Đáp án trả lời ngắn không ghi đơn vị.
}
\end{ex}

% ===== Câu 100 =====
% ID: AIQ_L10C3_FULL_1234
% Mức: VD
\begin{ex}
% ID: AIQ_L10C3_FULL_1234
% Mức: VD
Vật $4\,\text{kg}$ trên bàn ngang có ma sát $\mu=0,4$ nối qua ròng rọc với vật treo $6\,\text{kg}$. Bỏ qua khối lượng dây, ròng rọc. Lấy g= $10\,\text{m}$ /s^2. Tính gia tốc hệ.
\shortans{4,4}
\loigiai{
Với toàn hệ: $a=(m_2g-\mu m_1g)/(m_1+m_2)=4,4$ m/s^2. Đáp án trả lời ngắn không ghi đơn vị.
}
\end{ex}

% ===== Câu 101 =====
% ID: AIQ_L10C3_FULL_1235
% Mức: VD
\begin{ex}
% ID: AIQ_L10C3_FULL_1235
% Mức: VD
Vật $5\,\text{kg}$ trên bàn ngang có ma sát $\mu=0,1$ nối qua ròng rọc với vật treo $7\,\text{kg}$. Bỏ qua khối lượng dây, ròng rọc. Lấy g= $10\,\text{m}$ /s^2. Tính gia tốc hệ.
\shortans{5,42}
\loigiai{
Với toàn hệ: $a=(m_2g-\mu m_1g)/(m_1+m_2)=5,42$ m/s^2. Đáp án trả lời ngắn không ghi đơn vị.
}
\end{ex}

% ===== Câu 102 =====
% ID: AIQ_L10C3_FULL_1236
% Mức: VDC
\begin{ex}
% ID: AIQ_L10C3_FULL_1236
% Mức: VDC
Vật $6\,\text{kg}$ trên bàn ngang có ma sát $\mu=0,2$ nối qua ròng rọc với vật treo $8\,\text{kg}$. Bỏ qua khối lượng dây, ròng rọc. Lấy g= $10\,\text{m}$ /s^2. Tính gia tốc hệ.
\shortans{4,86}
\loigiai{
Với toàn hệ: $a=(m_2g-\mu m_1g)/(m_1+m_2)=4,86$ m/s^2. Đáp án trả lời ngắn không ghi đơn vị.
}
\end{ex}

% ===== Câu 103 =====
% ID: AIQ_L10C3_FULL_1237
% Mức: TH
\begin{bt}
% ID: AIQ_L10C3_FULL_1237
% Mức: TH
Vật $7\,\text{kg}$ trên bàn ngang có ma sát $\mu=0,25$ nối qua ròng rọc với vật treo $9\,\text{kg}$. Bỏ qua khối lượng dây, ròng rọc. Lấy g= $10\,\text{m}$ /s^2. Tính gia tốc hệ.
\loigiai{
Phân tích dữ kiện, chọn định luật phù hợp. Với toàn hệ: $a=(m_2g-\mu m_1g)/(m_1+m_2)=4,53$ m/s^2.
}
\end{bt}

% ===== Câu 104 =====
% ID: AIQ_L10C3_FULL_1238
% Mức: TH
\begin{bt}
% ID: AIQ_L10C3_FULL_1238
% Mức: TH
Vật $8\,\text{kg}$ trên bàn ngang có ma sát $\mu=0,3$ nối qua ròng rọc với vật treo $10\,\text{kg}$. Bỏ qua khối lượng dây, ròng rọc. Lấy g= $10\,\text{m}$ /s^2. Tính gia tốc hệ.
\loigiai{
Phân tích dữ kiện, chọn định luật phù hợp. Với toàn hệ: $a=(m_2g-\mu m_1g)/(m_1+m_2)=4,22$ m/s^2.
}
\end{bt}

% ===== Câu 105 =====
% ID: AIQ_L10C3_FULL_1239
% Mức: VD
\begin{bt}
% ID: AIQ_L10C3_FULL_1239
% Mức: VD
Vật $2\,\text{kg}$ trên bàn ngang có ma sát $\mu=0,4$ nối qua ròng rọc với vật treo $4\,\text{kg}$. Bỏ qua khối lượng dây, ròng rọc. Lấy g= $10\,\text{m}$ /s^2. Tính gia tốc hệ.
\loigiai{
Phân tích dữ kiện, chọn định luật phù hợp. Với toàn hệ: $a=(m_2g-\mu m_1g)/(m_1+m_2)=5,33$ m/s^2.
}
\end{bt}

% ===== Câu 106 =====
% ID: AIQ_L10C3_FULL_1240
% Mức: VD
\begin{bt}
% ID: AIQ_L10C3_FULL_1240
% Mức: VD
Vật $3\,\text{kg}$ trên bàn ngang có ma sát $\mu=0,1$ nối qua ròng rọc với vật treo $5\,\text{kg}$. Bỏ qua khối lượng dây, ròng rọc. Lấy g= $10\,\text{m}$ /s^2. Tính gia tốc hệ.
\loigiai{
Phân tích dữ kiện, chọn định luật phù hợp. Với toàn hệ: $a=(m_2g-\mu m_1g)/(m_1+m_2)=5,88$ m/s^2.
}
\end{bt}

% ===== Câu 107 =====
% ID: AIQ_L10C3_FULL_1241
% Mức: VD
\begin{bt}
% ID: AIQ_L10C3_FULL_1241
% Mức: VD
Vật $4\,\text{kg}$ trên bàn ngang có ma sát $\mu=0,2$ nối qua ròng rọc với vật treo $6\,\text{kg}$. Bỏ qua khối lượng dây, ròng rọc. Lấy g= $10\,\text{m}$ /s^2. Tính gia tốc hệ.
\loigiai{
Phân tích dữ kiện, chọn định luật phù hợp. Với toàn hệ: $a=(m_2g-\mu m_1g)/(m_1+m_2)=5,2$ m/s^2.
}
\end{bt}

% ===== Câu 108 =====
% ID: AIQ_L10C3_FULL_1242
% Mức: VDC
\begin{bt}
% ID: AIQ_L10C3_FULL_1242
% Mức: VDC
Vật $5\,\text{kg}$ trên bàn ngang có ma sát $\mu=0,25$ nối qua ròng rọc với vật treo $7\,\text{kg}$. Bỏ qua khối lượng dây, ròng rọc. Lấy g= $10\,\text{m}$ /s^2. Tính gia tốc hệ.
\loigiai{
Phân tích dữ kiện, chọn định luật phù hợp. Với toàn hệ: $a=(m_2g-\mu m_1g)/(m_1+m_2)=4,79$ m/s^2.
}
\end{bt}

% ===== Câu 109 =====
% ID: AIQ_L10C3_FULL_1243
% Mức: NB
\dangbt{Chuyển động của vật chịu lực kéo xiên}
\begin{ex}
% ID: AIQ_L10C3_FULL_1243
% Mức: NB
Kéo vật $6\,\text{kg}$ trên sàn bằng lực 29 $N$ hợp phương ngang góc 30°, $\mu=0,3$, g= $10\,\text{m}$ /s^2. Tính gia tốc.
\choice
{$3,91 m/s^2$}
{\True $1,91 m/s^2$}
{$2,29 m/s^2$}
{$1,53 m/s^2$}
\loigiai{
$N=mg-F\sin30^\circ=45,5N$; $f=\mu N=13,65N$; $a=(F\cos30^\circ-f)/m=1,91$ m/s^2. Chọn B.
}
\end{ex}

% ===== Câu 110 =====
% ID: AIQ_L10C3_FULL_1244
% Mức: NB
\begin{ex}
% ID: AIQ_L10C3_FULL_1244
% Mức: NB
Kéo vật $7\,\text{kg}$ trên sàn bằng lực 23 $N$ hợp phương ngang góc 30°, $\mu=0,4$, g= $10\,\text{m}$ /s^2. Tính gia tốc.
\choice
{\True $-0,5 m/s^2$}
{$-0,6 m/s^2$}
{$0 m/s^2$}
{$1,5 m/s^2$}
\loigiai{
$N=mg-F\sin30^\circ=58,5N$; $f=\mu N=23,4N$; $a=(F\cos30^\circ-f)/m=-0,5$ m/s^2. Chọn A.
}
\end{ex}

% ===== Câu 111 =====
% ID: AIQ_L10C3_FULL_1245
% Mức: NB
\begin{ex}
% ID: AIQ_L10C3_FULL_1245
% Mức: NB
Kéo vật $8\,\text{kg}$ trên sàn bằng lực 25 $N$ hợp phương ngang góc 30°, $\mu=0,1$, g= $10\,\text{m}$ /s^2. Tính gia tốc.
\choice
{$2,23 m/s^2$}
{$1,49 m/s^2$}
{$3,86 m/s^2$}
{\True $1,86 m/s^2$}
\loigiai{
$N=mg-F\sin30^\circ=67,5N$; $f=\mu N=6,75N$; $a=(F\cos30^\circ-f)/m=1,86$ m/s^2. Chọn D.
}
\end{ex}

% ===== Câu 112 =====
% ID: AIQ_L10C3_FULL_1246
% Mức: TH
\begin{ex}
% ID: AIQ_L10C3_FULL_1246
% Mức: TH
Kéo vật $2\,\text{kg}$ trên sàn bằng lực 27 $N$ hợp phương ngang góc 30°, $\mu=0,2$, g= $10\,\text{m}$ /s^2. Tính gia tốc.
\choice
{$8,83 m/s^2$}
{$13,04 m/s^2$}
{\True $11,04 m/s^2$}
{$13,25 m/s^2$}
\loigiai{
$N=mg-F\sin30^\circ=6,5N$; $f=\mu N=1,3N$; $a=(F\cos30^\circ-f)/m=11,04$ m/s^2. Chọn C.
}
\end{ex}

% ===== Câu 113 =====
% ID: AIQ_L10C3_FULL_1247
% Mức: TH
\begin{ex}
% ID: AIQ_L10C3_FULL_1247
% Mức: TH
Kéo vật $3\,\text{kg}$ trên sàn bằng lực 29 $N$ hợp phương ngang góc 30°, $\mu=0,25$, g= $10\,\text{m}$ /s^2. Tính gia tốc.
\choice
{$9,08 m/s^2$}
{\True $7,08 m/s^2$}
{$8,5 m/s^2$}
{$5,66 m/s^2$}
\loigiai{
$N=mg-F\sin30^\circ=15,5N$; $f=\mu N=3,88N$; $a=(F\cos30^\circ-f)/m=7,08$ m/s^2. Chọn B.
}
\end{ex}

% ===== Câu 114 =====
% ID: AIQ_L10C3_FULL_1248
% Mức: TH
\begin{ex}
% ID: AIQ_L10C3_FULL_1248
% Mức: TH
Kéo vật $4\,\text{kg}$ trên sàn bằng lực 23 $N$ hợp phương ngang góc 30°, $\mu=0,3$, g= $10\,\text{m}$ /s^2. Tính gia tốc.
\choice
{\True $2,84 m/s^2$}
{$3,41 m/s^2$}
{$2,27 m/s^2$}
{$4,84 m/s^2$}
\loigiai{
$N=mg-F\sin30^\circ=28,5N$; $f=\mu N=8,55N$; $a=(F\cos30^\circ-f)/m=2,84$ m/s^2. Chọn A.
}
\end{ex}

% ===== Câu 115 =====
% ID: AIQ_L10C3_FULL_1249
% Mức: TH
\begin{ex}
% ID: AIQ_L10C3_FULL_1249
% Mức: TH
Kéo vật $5\,\text{kg}$ trên sàn bằng lực 25 $N$ hợp phương ngang góc 30°, $\mu=0,4$, g= $10\,\text{m}$ /s^2. Tính gia tốc.
\choice
{$1,6 m/s^2$}
{$1,06 m/s^2$}
{$3,33 m/s^2$}
{\True $1,33 m/s^2$}
\loigiai{
$N=mg-F\sin30^\circ=37,5N$; $f=\mu N=15N$; $a=(F\cos30^\circ-f)/m=1,33$ m/s^2. Chọn D.
}
\end{ex}

% ===== Câu 116 =====
% ID: AIQ_L10C3_FULL_1250
% Mức: VD
\begin{ex}
% ID: AIQ_L10C3_FULL_1250
% Mức: VD
Kéo vật $6\,\text{kg}$ trên sàn bằng lực 27 $N$ hợp phương ngang góc 30°, $\mu=0,1$, g= $10\,\text{m}$ /s^2. Tính gia tốc.
\choice
{$2,5 m/s^2$}
{$5,12 m/s^2$}
{\True $3,12 m/s^2$}
{$3,74 m/s^2$}
\loigiai{
$N=mg-F\sin30^\circ=46,5N$; $f=\mu N=4,65N$; $a=(F\cos30^\circ-f)/m=3,12$ m/s^2. Chọn C.
}
\end{ex}

% ===== Câu 117 =====
% ID: AIQ_L10C3_FULL_1251
% Mức: VD
\begin{ex}
% ID: AIQ_L10C3_FULL_1251
% Mức: VD
Kéo vật $7\,\text{kg}$ trên sàn bằng lực 29 $N$ hợp phương ngang góc 30°, $\mu=0,2$, g= $10\,\text{m}$ /s^2. Tính gia tốc.
\choice
{$4 m/s^2$}
{\True $2 m/s^2$}
{$2,4 m/s^2$}
{$1,6 m/s^2$}
\loigiai{
$N=mg-F\sin30^\circ=55,5N$; $f=\mu N=11,1N$; $a=(F\cos30^\circ-f)/m=2$ m/s^2. Chọn B.
}
\end{ex}

% ===== Câu 118 =====
% ID: AIQ_L10C3_FULL_1252
% Mức: VDC
\begin{ex}
% ID: AIQ_L10C3_FULL_1252
% Mức: VDC
Kéo vật $8\,\text{kg}$ trên sàn bằng lực 23 $N$ hợp phương ngang góc 30°, $\mu=0,25$, g= $10\,\text{m}$ /s^2. Tính gia tốc.
\choice
{\True $0,35 m/s^2$}
{$0,42 m/s^2$}
{$0,28 m/s^2$}
{$2,35 m/s^2$}
\loigiai{
$N=mg-F\sin30^\circ=68,5N$; $f=\mu N=17,13N$; $a=(F\cos30^\circ-f)/m=0,35$ m/s^2. Chọn A.
}
\end{ex}

% ===== Câu 119 =====
% ID: AIQ_L10C3_FULL_1253
% Mức: NB
\begin{ex}
% ID: AIQ_L10C3_FULL_1253
% Mức: NB
Kéo vật $2\,\text{kg}$ trên sàn bằng lực 25 $N$ hợp phương ngang góc 30°, $\mu=0,3$, g= $10\,\text{m}$ /s^2. Tính gia tốc.
\choiceTF
{\True Kết quả đúng là $9,7\,\text{m}$ /s^2.}
{Kết quả bằng $19,4\,\text{m}$ /s^2.}
{\True Khi hợp lực bằng 0, vật có gia tốc bằng 0.}
{Có thể bỏ qua đơn vị của đại lượng trong kết quả vật lí.}
\loigiai{
\begin{itemchoice}
\item (Đúng) Kết quả đúng là $9,7\,\text{m}$ /s^2. (Vì): $N=mg-F\sin30^\circ=7,5N$; $f=\mu N=2,25N$; $a=(F\cos30^\circ-f)/m=9,7$ m/s^2. b) Sai vì giá trị hoặc chiều nêu ra không phù hợp. c) Đúng theo định luật II Newton. d) Sai vì kết quả vật lí phải kèm đơn vị thích hợp
\item (Sai) Kết quả bằng $19,4\,\text{m}$ /s^2.
\item (Đúng) Khi hợp lực bằng 0, vật có gia tốc bằng 0.
\item (Sai) Có thể bỏ qua đơn vị của đại lượng trong kết quả vật lí.
\end{itemchoice}
}
\end{ex}

% ===== Câu 120 =====
% ID: AIQ_L10C3_FULL_1254
% Mức: NB
\begin{ex}
% ID: AIQ_L10C3_FULL_1254
% Mức: NB
Kéo vật $3\,\text{kg}$ trên sàn bằng lực 27 $N$ hợp phương ngang góc 30°, $\mu=0,4$, g= $10\,\text{m}$ /s^2. Tính gia tốc.
\choiceTF
{\True Kết quả đúng là $5,59\,\text{m}$ /s^2.}
{Kết quả bằng $11,18\,\text{m}$ /s^2.}
{\True Khi hợp lực bằng 0, vật có gia tốc bằng 0.}
{Có thể bỏ qua đơn vị của đại lượng trong kết quả vật lí.}
\loigiai{
\begin{itemchoice}
\item (Đúng) Kết quả đúng là $5,59\,\text{m}$ /s^2. (Vì): $N=mg-F\sin30^\circ=16,5N$; $f=\mu N=6,6N$; $a=(F\cos30^\circ-f)/m=5,59$ m/s^2. b) Sai vì giá trị hoặc chiều nêu ra không phù hợp. c) Đúng theo định luật II Newton. d) Sai vì kết quả vật lí phải kèm đơn vị thích hợp
\item (Sai) Kết quả bằng $11,18\,\text{m}$ /s^2.
\item (Đúng) Khi hợp lực bằng 0, vật có gia tốc bằng 0.
\item (Sai) Có thể bỏ qua đơn vị của đại lượng trong kết quả vật lí.
\end{itemchoice}
}
\end{ex}

% ===== Câu 121 =====
% ID: AIQ_L10C3_FULL_1255
% Mức: TH
\begin{ex}
% ID: AIQ_L10C3_FULL_1255
% Mức: TH
Kéo vật $4\,\text{kg}$ trên sàn bằng lực 29 $N$ hợp phương ngang góc 30°, $\mu=0,1$, g= $10\,\text{m}$ /s^2. Tính gia tốc.
\choiceTF
{\True Kết quả đúng là $5,64\,\text{m}$ /s^2.}
{Kết quả bằng $11,28\,\text{m}$ /s^2.}
{\True Khi hợp lực bằng 0, vật có gia tốc bằng 0.}
{Có thể bỏ qua đơn vị của đại lượng trong kết quả vật lí.}
\loigiai{
\begin{itemchoice}
\item (Đúng) Kết quả đúng là $5,64\,\text{m}$ /s^2. (Vì): $N=mg-F\sin30^\circ=25,5N$; $f=\mu N=2,55N$; $a=(F\cos30^\circ-f)/m=5,64$ m/s^2. b) Sai vì giá trị hoặc chiều nêu ra không phù hợp. c) Đúng theo định luật II Newton. d) Sai vì kết quả vật lí phải kèm đơn vị thích hợp
\item (Sai) Kết quả bằng $11,28\,\text{m}$ /s^2.
\item (Đúng) Khi hợp lực bằng 0, vật có gia tốc bằng 0.
\item (Sai) Có thể bỏ qua đơn vị của đại lượng trong kết quả vật lí.
\end{itemchoice}
}
\end{ex}

% ===== Câu 122 =====
% ID: AIQ_L10C3_FULL_1256
% Mức: TH
\begin{ex}
% ID: AIQ_L10C3_FULL_1256
% Mức: TH
Kéo vật $5\,\text{kg}$ trên sàn bằng lực 23 $N$ hợp phương ngang góc 30°, $\mu=0,2$, g= $10\,\text{m}$ /s^2. Tính gia tốc.
\choiceTF
{\True Kết quả đúng là $2,44\,\text{m}$ /s^2.}
{Kết quả bằng $4,88\,\text{m}$ /s^2.}
{\True Khi hợp lực bằng 0, vật có gia tốc bằng 0.}
{Có thể bỏ qua đơn vị của đại lượng trong kết quả vật lí.}
\loigiai{
\begin{itemchoice}
\item (Đúng) Kết quả đúng là $2,44\,\text{m}$ /s^2. (Vì): $N=mg-F\sin30^\circ=38,5N$; $f=\mu N=7,7N$; $a=(F\cos30^\circ-f)/m=2,44$ m/s^2. b) Sai vì giá trị hoặc chiều nêu ra không phù hợp. c) Đúng theo định luật II Newton. d) Sai vì kết quả vật lí phải kèm đơn vị thích hợp
\item (Sai) Kết quả bằng $4,88\,\text{m}$ /s^2.
\item (Đúng) Khi hợp lực bằng 0, vật có gia tốc bằng 0.
\item (Sai) Có thể bỏ qua đơn vị của đại lượng trong kết quả vật lí.
\end{itemchoice}
}
\end{ex}

% ===== Câu 123 =====
% ID: AIQ_L10C3_FULL_1257
% Mức: VD
\begin{ex}
% ID: AIQ_L10C3_FULL_1257
% Mức: VD
Kéo vật $6\,\text{kg}$ trên sàn bằng lực 25 $N$ hợp phương ngang góc 30°, $\mu=0,25$, g= $10\,\text{m}$ /s^2. Tính gia tốc.
\choiceTF
{\True Kết quả đúng là $1,63\,\text{m}$ /s^2.}
{Kết quả bằng $3,26\,\text{m}$ /s^2.}
{\True Khi hợp lực bằng 0, vật có gia tốc bằng 0.}
{Có thể bỏ qua đơn vị của đại lượng trong kết quả vật lí.}
\loigiai{
\begin{itemchoice}
\item (Đúng) Kết quả đúng là $1,63\,\text{m}$ /s^2. (Vì): $N=mg-F\sin30^\circ=47,5N$; $f=\mu N=11,88N$; $a=(F\cos30^\circ-f)/m=1,63$ m/s^2. b) Sai vì giá trị hoặc chiều nêu ra không phù hợp. c) Đúng theo định luật II Newton. d) Sai vì kết quả vật lí phải kèm đơn vị thích hợp
\item (Sai) Kết quả bằng $3,26\,\text{m}$ /s^2.
\item (Đúng) Khi hợp lực bằng 0, vật có gia tốc bằng 0.
\item (Sai) Có thể bỏ qua đơn vị của đại lượng trong kết quả vật lí.
\end{itemchoice}
}
\end{ex}

% ===== Câu 124 =====
% ID: AIQ_L10C3_FULL_1258
% Mức: TH
\begin{ex}
% ID: AIQ_L10C3_FULL_1258
% Mức: TH
Kéo vật $7\,\text{kg}$ trên sàn bằng lực 27 $N$ hợp phương ngang góc 30°, $\mu=0,3$, g= $10\,\text{m}$ /s^2. Tính gia tốc.
\shortans{0,92}
\loigiai{
$N=mg-F\sin30^\circ=56,5N$; $f=\mu N=16,95N$; $a=(F\cos30^\circ-f)/m=0,92$ m/s^2. Đáp án trả lời ngắn không ghi đơn vị.
}
\end{ex}

% ===== Câu 125 =====
% ID: AIQ_L10C3_FULL_1259
% Mức: TH
\begin{ex}
% ID: AIQ_L10C3_FULL_1259
% Mức: TH
Kéo vật $8\,\text{kg}$ trên sàn bằng lực 29 $N$ hợp phương ngang góc 30°, $\mu=0,4$, g= $10\,\text{m}$ /s^2. Tính gia tốc.
\shortans{-0,14}
\loigiai{
$N=mg-F\sin30^\circ=65,5N$; $f=\mu N=26,2N$; $a=(F\cos30^\circ-f)/m=-0,14$ m/s^2. Đáp án trả lời ngắn không ghi đơn vị.
}
\end{ex}

% ===== Câu 126 =====
% ID: AIQ_L10C3_FULL_1260
% Mức: VD
\begin{ex}
% ID: AIQ_L10C3_FULL_1260
% Mức: VD
Kéo vật $2\,\text{kg}$ trên sàn bằng lực 23 $N$ hợp phương ngang góc 30°, $\mu=0,1$, g= $10\,\text{m}$ /s^2. Tính gia tốc.
\shortans{9,53}
\loigiai{
$N=mg-F\sin30^\circ=8,5N$; $f=\mu N=0,85N$; $a=(F\cos30^\circ-f)/m=9,53$ m/s^2. Đáp án trả lời ngắn không ghi đơn vị.
}
\end{ex}

% ===== Câu 127 =====
% ID: AIQ_L10C3_FULL_1261
% Mức: VD
\begin{ex}
% ID: AIQ_L10C3_FULL_1261
% Mức: VD
Kéo vật $3\,\text{kg}$ trên sàn bằng lực 25 $N$ hợp phương ngang góc 30°, $\mu=0,2$, g= $10\,\text{m}$ /s^2. Tính gia tốc.
\shortans{6,05}
\loigiai{
$N=mg-F\sin30^\circ=17,5N$; $f=\mu N=3,5N$; $a=(F\cos30^\circ-f)/m=6,05$ m/s^2. Đáp án trả lời ngắn không ghi đơn vị.
}
\end{ex}

% ===== Câu 128 =====
% ID: AIQ_L10C3_FULL_1262
% Mức: VD
\begin{ex}
% ID: AIQ_L10C3_FULL_1262
% Mức: VD
Kéo vật $4\,\text{kg}$ trên sàn bằng lực 27 $N$ hợp phương ngang góc 30°, $\mu=0,25$, g= $10\,\text{m}$ /s^2. Tính gia tốc.
\shortans{4,19}
\loigiai{
$N=mg-F\sin30^\circ=26,5N$; $f=\mu N=6,63N$; $a=(F\cos30^\circ-f)/m=4,19$ m/s^2. Đáp án trả lời ngắn không ghi đơn vị.
}
\end{ex}

% ===== Câu 129 =====
% ID: AIQ_L10C3_FULL_1263
% Mức: VDC
\begin{ex}
% ID: AIQ_L10C3_FULL_1263
% Mức: VDC
Kéo vật $5\,\text{kg}$ trên sàn bằng lực 29 $N$ hợp phương ngang góc 30°, $\mu=0,3$, g= $10\,\text{m}$ /s^2. Tính gia tốc.
\shortans{2,89}
\loigiai{
$N=mg-F\sin30^\circ=35,5N$; $f=\mu N=10,65N$; $a=(F\cos30^\circ-f)/m=2,89$ m/s^2. Đáp án trả lời ngắn không ghi đơn vị.
}
\end{ex}

% ===== Câu 130 =====
% ID: AIQ_L10C3_FULL_1264
% Mức: TH
\begin{bt}
% ID: AIQ_L10C3_FULL_1264
% Mức: TH
Kéo vật $6\,\text{kg}$ trên sàn bằng lực 23 $N$ hợp phương ngang góc 30°, $\mu=0,4$, g= $10\,\text{m}$ /s^2. Tính gia tốc.
\loigiai{
Phân tích dữ kiện, chọn định luật phù hợp. $N=mg-F\sin30^\circ=48,5N$; $f=\mu N=19,4N$; $a=(F\cos30^\circ-f)/m=0,09$ m/s^2.
}
\end{bt}

% ===== Câu 131 =====
% ID: AIQ_L10C3_FULL_1265
% Mức: TH
\begin{bt}
% ID: AIQ_L10C3_FULL_1265
% Mức: TH
Kéo vật $7\,\text{kg}$ trên sàn bằng lực 25 $N$ hợp phương ngang góc 30°, $\mu=0,1$, g= $10\,\text{m}$ /s^2. Tính gia tốc.
\loigiai{
Phân tích dữ kiện, chọn định luật phù hợp. $N=mg-F\sin30^\circ=57,5N$; $f=\mu N=5,75N$; $a=(F\cos30^\circ-f)/m=2,27$ m/s^2.
}
\end{bt}

% ===== Câu 132 =====
% ID: AIQ_L10C3_FULL_1266
% Mức: VD
\begin{bt}
% ID: AIQ_L10C3_FULL_1266
% Mức: VD
Kéo vật $8\,\text{kg}$ trên sàn bằng lực 27 $N$ hợp phương ngang góc 30°, $\mu=0,2$, g= $10\,\text{m}$ /s^2. Tính gia tốc.
\loigiai{
Phân tích dữ kiện, chọn định luật phù hợp. $N=mg-F\sin30^\circ=66,5N$; $f=\mu N=13,3N$; $a=(F\cos30^\circ-f)/m=1,26$ m/s^2.
}
\end{bt}

% ===== Câu 133 =====
% ID: AIQ_L10C3_FULL_1267
% Mức: VD
\begin{bt}
% ID: AIQ_L10C3_FULL_1267
% Mức: VD
Kéo vật $2\,\text{kg}$ trên sàn bằng lực 29 $N$ hợp phương ngang góc 30°, $\mu=0,25$, g= $10\,\text{m}$ /s^2. Tính gia tốc.
\loigiai{
Phân tích dữ kiện, chọn định luật phù hợp. $N=mg-F\sin30^\circ=5,5N$; $f=\mu N=1,38N$; $a=(F\cos30^\circ-f)/m=11,87$ m/s^2.
}
\end{bt}

% ===== Câu 134 =====
% ID: AIQ_L10C3_FULL_1268
% Mức: VD
\begin{bt}
% ID: AIQ_L10C3_FULL_1268
% Mức: VD
Kéo vật $3\,\text{kg}$ trên sàn bằng lực 23 $N$ hợp phương ngang góc 30°, $\mu=0,3$, g= $10\,\text{m}$ /s^2. Tính gia tốc.
\loigiai{
Phân tích dữ kiện, chọn định luật phù hợp. $N=mg-F\sin30^\circ=18,5N$; $f=\mu N=5,55N$; $a=(F\cos30^\circ-f)/m=4,79$ m/s^2.
}
\end{bt}

% ===== Câu 135 =====
% ID: AIQ_L10C3_FULL_1269
% Mức: VDC
\begin{bt}
% ID: AIQ_L10C3_FULL_1269
% Mức: VDC
Kéo vật $4\,\text{kg}$ trên sàn bằng lực 25 $N$ hợp phương ngang góc 30°, $\mu=0,4$, g= $10\,\text{m}$ /s^2. Tính gia tốc.
\loigiai{
Phân tích dữ kiện, chọn định luật phù hợp. $N=mg-F\sin30^\circ=27,5N$; $f=\mu N=11N$; $a=(F\cos30^\circ-f)/m=2,66$ m/s^2.
}
\end{bt}

% ===== Câu 136 =====
% ID: AIQ_L10C3_FULL_1270
% Mức: NB
\dangbt{Bài toán động lực học tổng hợp và ứng dụng thực tiễn}
\begin{ex}
% ID: AIQ_L10C3_FULL_1270
% Mức: NB
Xe đồ chơi khối lượng $5\,\text{kg}$ bắt đầu từ nghỉ, chịu hợp lực không đổi 32 $N$ trong $4\,\text{s}$. Bỏ qua lực cản. Tính quãng đường.
\choice
{$40,96\,\text{m}$}
{$53,2\,\text{m}$}
{\True $51,2\,\text{m}$}
{$61,44\,\text{m}$}
\loigiai{
$a=F/m=6,4m/s^2;s=\frac12at^2=51,2$ m. Chọn C.
}
\end{ex}

% ===== Câu 137 =====
% ID: AIQ_L10C3_FULL_1271
% Mức: NB
\begin{ex}
% ID: AIQ_L10C3_FULL_1271
% Mức: NB
Xe đồ chơi khối lượng $6\,\text{kg}$ bắt đầu từ nghỉ, chịu hợp lực không đổi 34 $N$ trong $5\,\text{s}$. Bỏ qua lực cản. Tính quãng đường.
\choice
{$72,88\,\text{m}$}
{\True $70,88\,\text{m}$}
{$85,06\,\text{m}$}
{$56,7\,\text{m}$}
\loigiai{
$a=F/m=5,67m/s^2;s=\frac12at^2=70,88$ m. Chọn B.
}
\end{ex}

% ===== Câu 138 =====
% ID: AIQ_L10C3_FULL_1272
% Mức: NB
\begin{ex}
% ID: AIQ_L10C3_FULL_1272
% Mức: NB
Xe đồ chơi khối lượng $7\,\text{kg}$ bắt đầu từ nghỉ, chịu hợp lực không đổi 28 $N$ trong $2\,\text{s}$. Bỏ qua lực cản. Tính quãng đường.
\choice
{\True $8\,\text{m}$}
{$9,6\,\text{m}$}
{$6,4\,\text{m}$}
{$10\,\text{m}$}
\loigiai{
$a=F/m=4m/s^2;s=\frac12at^2=8$ m. Chọn A.
}
\end{ex}

% ===== Câu 139 =====
% ID: AIQ_L10C3_FULL_1273
% Mức: TH
\begin{ex}
% ID: AIQ_L10C3_FULL_1273
% Mức: TH
Xe đồ chơi khối lượng $8\,\text{kg}$ bắt đầu từ nghỉ, chịu hợp lực không đổi 30 $N$ trong $3\,\text{s}$. Bỏ qua lực cản. Tính quãng đường.
\choice
{$20,26\,\text{m}$}
{$13,5\,\text{m}$}
{$18,88\,\text{m}$}
{\True $16,88\,\text{m}$}
\loigiai{
$a=F/m=3,75m/s^2;s=\frac12at^2=16,88$ m. Chọn D.
}
\end{ex}

% ===== Câu 140 =====
% ID: AIQ_L10C3_FULL_1274
% Mức: TH
\begin{ex}
% ID: AIQ_L10C3_FULL_1274
% Mức: TH
Xe đồ chơi khối lượng $2\,\text{kg}$ bắt đầu từ nghỉ, chịu hợp lực không đổi 32 $N$ trong $4\,\text{s}$. Bỏ qua lực cản. Tính quãng đường.
\choice
{$102,4\,\text{m}$}
{$130\,\text{m}$}
{\True $128\,\text{m}$}
{$153,6\,\text{m}$}
\loigiai{
$a=F/m=16m/s^2;s=\frac12at^2=128$ m. Chọn C.
}
\end{ex}

% ===== Câu 141 =====
% ID: AIQ_L10C3_FULL_1275
% Mức: TH
\begin{ex}
% ID: AIQ_L10C3_FULL_1275
% Mức: TH
Xe đồ chơi khối lượng $3\,\text{kg}$ bắt đầu từ nghỉ, chịu hợp lực không đổi 34 $N$ trong $5\,\text{s}$. Bỏ qua lực cản. Tính quãng đường.
\choice
{$143,63\,\text{m}$}
{\True $141,63\,\text{m}$}
{$169,96\,\text{m}$}
{$113,3\,\text{m}$}
\loigiai{
$a=F/m=11,33m/s^2;s=\frac12at^2=141,63$ m. Chọn B.
}
\end{ex}

% ===== Câu 142 =====
% ID: AIQ_L10C3_FULL_1276
% Mức: TH
\begin{ex}
% ID: AIQ_L10C3_FULL_1276
% Mức: TH
Xe đồ chơi khối lượng $4\,\text{kg}$ bắt đầu từ nghỉ, chịu hợp lực không đổi 28 $N$ trong $2\,\text{s}$. Bỏ qua lực cản. Tính quãng đường.
\choice
{\True $14\,\text{m}$}
{$16,8\,\text{m}$}
{$11,2\,\text{m}$}
{$16\,\text{m}$}
\loigiai{
$a=F/m=7m/s^2;s=\frac12at^2=14$ m. Chọn A.
}
\end{ex}

% ===== Câu 143 =====
% ID: AIQ_L10C3_FULL_1277
% Mức: VD
\begin{ex}
% ID: AIQ_L10C3_FULL_1277
% Mức: VD
Xe đồ chơi khối lượng $5\,\text{kg}$ bắt đầu từ nghỉ, chịu hợp lực không đổi 30 $N$ trong $3\,\text{s}$. Bỏ qua lực cản. Tính quãng đường.
\choice
{$32,4\,\text{m}$}
{$21,6\,\text{m}$}
{$29\,\text{m}$}
{\True $27\,\text{m}$}
\loigiai{
$a=F/m=6m/s^2;s=\frac12at^2=27$ m. Chọn D.
}
\end{ex}

% ===== Câu 144 =====
% ID: AIQ_L10C3_FULL_1278
% Mức: VD
\begin{ex}
% ID: AIQ_L10C3_FULL_1278
% Mức: VD
Xe đồ chơi khối lượng $6\,\text{kg}$ bắt đầu từ nghỉ, chịu hợp lực không đổi 32 $N$ trong $4\,\text{s}$. Bỏ qua lực cản. Tính quãng đường.
\choice
{$34,11\,\text{m}$}
{$44,64\,\text{m}$}
{\True $42,64\,\text{m}$}
{$51,17\,\text{m}$}
\loigiai{
$a=F/m=5,33m/s^2;s=\frac12at^2=42,64$ m. Chọn C.
}
\end{ex}

% ===== Câu 145 =====
% ID: AIQ_L10C3_FULL_1279
% Mức: VDC
\begin{ex}
% ID: AIQ_L10C3_FULL_1279
% Mức: VDC
Xe đồ chơi khối lượng $7\,\text{kg}$ bắt đầu từ nghỉ, chịu hợp lực không đổi 34 $N$ trong $5\,\text{s}$. Bỏ qua lực cản. Tính quãng đường.
\choice
{$62,75\,\text{m}$}
{\True $60,75\,\text{m}$}
{$72,9\,\text{m}$}
{$48,6\,\text{m}$}
\loigiai{
$a=F/m=4,86m/s^2;s=\frac12at^2=60,75$ m. Chọn B.
}
\end{ex}

% ===== Câu 146 =====
% ID: AIQ_L10C3_FULL_1280
% Mức: NB
\begin{ex}
% ID: AIQ_L10C3_FULL_1280
% Mức: NB
Xe đồ chơi khối lượng $8\,\text{kg}$ bắt đầu từ nghỉ, chịu hợp lực không đổi 28 $N$ trong $2\,\text{s}$. Bỏ qua lực cản. Tính quãng đường.
\choiceTF
{\True Kết quả đúng là $7\,\text{m}$.}
{Kết quả bằng $14\,\text{m}$.}
{\True Khi hợp lực bằng 0, vật có gia tốc bằng 0.}
{Có thể bỏ qua đơn vị của đại lượng trong kết quả vật lí.}
\loigiai{
\begin{itemchoice}
\item (Đúng) Kết quả đúng là $7\,\text{m}$. (Vì): $a=F/m=3,5m/s^2;s=\frac12at^2=7$ m. b) Sai vì giá trị hoặc chiều nêu ra không phù hợp. c) Đúng theo định luật II Newton. d) Sai vì kết quả vật lí phải kèm đơn vị thích hợp
\item (Sai) Kết quả bằng $14\,\text{m}$.
\item (Đúng) Khi hợp lực bằng 0, vật có gia tốc bằng 0.
\item (Sai) Có thể bỏ qua đơn vị của đại lượng trong kết quả vật lí.
\end{itemchoice}
}
\end{ex}

% ===== Câu 147 =====
% ID: AIQ_L10C3_FULL_1281
% Mức: NB
\begin{ex}
% ID: AIQ_L10C3_FULL_1281
% Mức: NB
Xe đồ chơi khối lượng $2\,\text{kg}$ bắt đầu từ nghỉ, chịu hợp lực không đổi 30 $N$ trong $3\,\text{s}$. Bỏ qua lực cản. Tính quãng đường.
\choiceTF
{\True Kết quả đúng là $67,5\,\text{m}$.}
{Kết quả bằng $135\,\text{m}$.}
{\True Khi hợp lực bằng 0, vật có gia tốc bằng 0.}
{Có thể bỏ qua đơn vị của đại lượng trong kết quả vật lí.}
\loigiai{
\begin{itemchoice}
\item (Đúng) Kết quả đúng là $67,5\,\text{m}$. (Vì): $a=F/m=15m/s^2;s=\frac12at^2=67,5$ m. b) Sai vì giá trị hoặc chiều nêu ra không phù hợp. c) Đúng theo định luật II Newton. d) Sai vì kết quả vật lí phải kèm đơn vị thích hợp
\item (Sai) Kết quả bằng $135\,\text{m}$.
\item (Đúng) Khi hợp lực bằng 0, vật có gia tốc bằng 0.
\item (Sai) Có thể bỏ qua đơn vị của đại lượng trong kết quả vật lí.
\end{itemchoice}
}
\end{ex}

% ===== Câu 148 =====
% ID: AIQ_L10C3_FULL_1282
% Mức: TH
\begin{ex}
% ID: AIQ_L10C3_FULL_1282
% Mức: TH
Xe đồ chơi khối lượng $3\,\text{kg}$ bắt đầu từ nghỉ, chịu hợp lực không đổi 32 $N$ trong $4\,\text{s}$. Bỏ qua lực cản. Tính quãng đường.
\choiceTF
{\True Kết quả đúng là $85,36\,\text{m}$.}
{Kết quả bằng $170,72\,\text{m}$.}
{\True Khi hợp lực bằng 0, vật có gia tốc bằng 0.}
{Có thể bỏ qua đơn vị của đại lượng trong kết quả vật lí.}
\loigiai{
\begin{itemchoice}
\item (Đúng) Kết quả đúng là $85,36\,\text{m}$. (Vì): $a=F/m=10,67m/s^2;s=\frac12at^2=85,36$ m. b) Sai vì giá trị hoặc chiều nêu ra không phù hợp. c) Đúng theo định luật II Newton. d) Sai vì kết quả vật lí phải kèm đơn vị thích hợp
\item (Sai) Kết quả bằng $170,72\,\text{m}$.
\item (Đúng) Khi hợp lực bằng 0, vật có gia tốc bằng 0.
\item (Sai) Có thể bỏ qua đơn vị của đại lượng trong kết quả vật lí.
\end{itemchoice}
}
\end{ex}

% ===== Câu 149 =====
% ID: AIQ_L10C3_FULL_1283
% Mức: TH
\begin{ex}
% ID: AIQ_L10C3_FULL_1283
% Mức: TH
Xe đồ chơi khối lượng $4\,\text{kg}$ bắt đầu từ nghỉ, chịu hợp lực không đổi 34 $N$ trong $5\,\text{s}$. Bỏ qua lực cản. Tính quãng đường.
\choiceTF
{\True Kết quả đúng là $106,25\,\text{m}$.}
{Kết quả bằng $212,5\,\text{m}$.}
{\True Khi hợp lực bằng 0, vật có gia tốc bằng 0.}
{Có thể bỏ qua đơn vị của đại lượng trong kết quả vật lí.}
\loigiai{
\begin{itemchoice}
\item (Đúng) Kết quả đúng là $106,25\,\text{m}$. (Vì): $a=F/m=8,5m/s^2;s=\frac12at^2=106,25$ m. b) Sai vì giá trị hoặc chiều nêu ra không phù hợp. c) Đúng theo định luật II Newton. d) Sai vì kết quả vật lí phải kèm đơn vị thích hợp
\item (Sai) Kết quả bằng $212,5\,\text{m}$.
\item (Đúng) Khi hợp lực bằng 0, vật có gia tốc bằng 0.
\item (Sai) Có thể bỏ qua đơn vị của đại lượng trong kết quả vật lí.
\end{itemchoice}
}
\end{ex}

% ===== Câu 150 =====
% ID: AIQ_L10C3_FULL_1284
% Mức: VD
\begin{ex}
% ID: AIQ_L10C3_FULL_1284
% Mức: VD
Xe đồ chơi khối lượng $5\,\text{kg}$ bắt đầu từ nghỉ, chịu hợp lực không đổi 28 $N$ trong $2\,\text{s}$. Bỏ qua lực cản. Tính quãng đường.
\choiceTF
{\True Kết quả đúng là $11,2\,\text{m}$.}
{Kết quả bằng $22,4\,\text{m}$.}
{\True Khi hợp lực bằng 0, vật có gia tốc bằng 0.}
{Có thể bỏ qua đơn vị của đại lượng trong kết quả vật lí.}
\loigiai{
\begin{itemchoice}
\item (Đúng) Kết quả đúng là $11,2\,\text{m}$. (Vì): $a=F/m=5,6m/s^2;s=\frac12at^2=11,2$ m. b) Sai vì giá trị hoặc chiều nêu ra không phù hợp. c) Đúng theo định luật II Newton. d) Sai vì kết quả vật lí phải kèm đơn vị thích hợp
\item (Sai) Kết quả bằng $22,4\,\text{m}$.
\item (Đúng) Khi hợp lực bằng 0, vật có gia tốc bằng 0.
\item (Sai) Có thể bỏ qua đơn vị của đại lượng trong kết quả vật lí.
\end{itemchoice}
}
\end{ex}

% ===== Câu 151 =====
% ID: AIQ_L10C3_FULL_1285
% Mức: TH
\begin{ex}
% ID: AIQ_L10C3_FULL_1285
% Mức: TH
Xe đồ chơi khối lượng $6\,\text{kg}$ bắt đầu từ nghỉ, chịu hợp lực không đổi 30 $N$ trong $3\,\text{s}$. Bỏ qua lực cản. Tính quãng đường.
\shortans{22,5}
\loigiai{
$a=F/m=5m/s^2;s=\frac12at^2=22,5$ m. Đáp án trả lời ngắn không ghi đơn vị.
}
\end{ex}

% ===== Câu 152 =====
% ID: AIQ_L10C3_FULL_1286
% Mức: TH
\begin{ex}
% ID: AIQ_L10C3_FULL_1286
% Mức: TH
Xe đồ chơi khối lượng $7\,\text{kg}$ bắt đầu từ nghỉ, chịu hợp lực không đổi 32 $N$ trong $4\,\text{s}$. Bỏ qua lực cản. Tính quãng đường.
\shortans{36,56}
\loigiai{
$a=F/m=4,57m/s^2;s=\frac12at^2=36,56$ m. Đáp án trả lời ngắn không ghi đơn vị.
}
\end{ex}

% ===== Câu 153 =====
% ID: AIQ_L10C3_FULL_1287
% Mức: VD
\begin{ex}
% ID: AIQ_L10C3_FULL_1287
% Mức: VD
Xe đồ chơi khối lượng $8\,\text{kg}$ bắt đầu từ nghỉ, chịu hợp lực không đổi 34 $N$ trong $5\,\text{s}$. Bỏ qua lực cản. Tính quãng đường.
\shortans{53,13}
\loigiai{
$a=F/m=4,25m/s^2;s=\frac12at^2=53,13$ m. Đáp án trả lời ngắn không ghi đơn vị.
}
\end{ex}

% ===== Câu 154 =====
% ID: AIQ_L10C3_FULL_1288
% Mức: VD
\begin{ex}
% ID: AIQ_L10C3_FULL_1288
% Mức: VD
Xe đồ chơi khối lượng $2\,\text{kg}$ bắt đầu từ nghỉ, chịu hợp lực không đổi 28 $N$ trong $2\,\text{s}$. Bỏ qua lực cản. Tính quãng đường.
\shortans{28}
\loigiai{
$a=F/m=14m/s^2;s=\frac12at^2=28$ m. Đáp án trả lời ngắn không ghi đơn vị.
}
\end{ex}

% ===== Câu 155 =====
% ID: AIQ_L10C3_FULL_1289
% Mức: VD
\begin{ex}
% ID: AIQ_L10C3_FULL_1289
% Mức: VD
Xe đồ chơi khối lượng $3\,\text{kg}$ bắt đầu từ nghỉ, chịu hợp lực không đổi 30 $N$ trong $3\,\text{s}$. Bỏ qua lực cản. Tính quãng đường.
\shortans{45}
\loigiai{
$a=F/m=10m/s^2;s=\frac12at^2=45$ m. Đáp án trả lời ngắn không ghi đơn vị.
}
\end{ex}

% ===== Câu 156 =====
% ID: AIQ_L10C3_FULL_1290
% Mức: VDC
\begin{ex}
% ID: AIQ_L10C3_FULL_1290
% Mức: VDC
Xe đồ chơi khối lượng $4\,\text{kg}$ bắt đầu từ nghỉ, chịu hợp lực không đổi 32 $N$ trong $4\,\text{s}$. Bỏ qua lực cản. Tính quãng đường.
\shortans{64}
\loigiai{
$a=F/m=8m/s^2;s=\frac12at^2=64$ m. Đáp án trả lời ngắn không ghi đơn vị.
}
\end{ex}

% ===== Câu 157 =====
% ID: AIQ_L10C3_FULL_1291
% Mức: TH
\begin{bt}
% ID: AIQ_L10C3_FULL_1291
% Mức: TH
Xe đồ chơi khối lượng $5\,\text{kg}$ bắt đầu từ nghỉ, chịu hợp lực không đổi 34 $N$ trong $5\,\text{s}$. Bỏ qua lực cản. Tính quãng đường.
\loigiai{
Phân tích dữ kiện, chọn định luật phù hợp. $a=F/m=6,8m/s^2;s=\frac12at^2=85$ m.
}
\end{bt}

% ===== Câu 158 =====
% ID: AIQ_L10C3_FULL_1292
% Mức: TH
\begin{bt}
% ID: AIQ_L10C3_FULL_1292
% Mức: TH
Xe đồ chơi khối lượng $6\,\text{kg}$ bắt đầu từ nghỉ, chịu hợp lực không đổi 28 $N$ trong $2\,\text{s}$. Bỏ qua lực cản. Tính quãng đường.
\loigiai{
Phân tích dữ kiện, chọn định luật phù hợp. $a=F/m=4,67m/s^2;s=\frac12at^2=9,34$ m.
}
\end{bt}

% ===== Câu 159 =====
% ID: AIQ_L10C3_FULL_1293
% Mức: VD
\begin{bt}
% ID: AIQ_L10C3_FULL_1293
% Mức: VD
Xe đồ chơi khối lượng $7\,\text{kg}$ bắt đầu từ nghỉ, chịu hợp lực không đổi 30 $N$ trong $3\,\text{s}$. Bỏ qua lực cản. Tính quãng đường.
\loigiai{
Phân tích dữ kiện, chọn định luật phù hợp. $a=F/m=4,29m/s^2;s=\frac12at^2=19,31$ m.
}
\end{bt}

% ===== Câu 160 =====
% ID: AIQ_L10C3_FULL_1294
% Mức: VD
\begin{bt}
% ID: AIQ_L10C3_FULL_1294
% Mức: VD
Xe đồ chơi khối lượng $8\,\text{kg}$ bắt đầu từ nghỉ, chịu hợp lực không đổi 32 $N$ trong $4\,\text{s}$. Bỏ qua lực cản. Tính quãng đường.
\loigiai{
Phân tích dữ kiện, chọn định luật phù hợp. $a=F/m=4m/s^2;s=\frac12at^2=32$ m.
}
\end{bt}

% ===== Câu 161 =====
% ID: AIQ_L10C3_FULL_1295
% Mức: VD
\begin{bt}
% ID: AIQ_L10C3_FULL_1295
% Mức: VD
Xe đồ chơi khối lượng $2\,\text{kg}$ bắt đầu từ nghỉ, chịu hợp lực không đổi 34 $N$ trong $5\,\text{s}$. Bỏ qua lực cản. Tính quãng đường.
\loigiai{
Phân tích dữ kiện, chọn định luật phù hợp. $a=F/m=17m/s^2;s=\frac12at^2=212,5$ m.
}
\end{bt}

% ===== Câu 162 =====
% ID: AIQ_L10C3_FULL_1296
% Mức: VDC
\begin{bt}
% ID: AIQ_L10C3_FULL_1296
% Mức: VDC
Xe đồ chơi khối lượng $3\,\text{kg}$ bắt đầu từ nghỉ, chịu hợp lực không đổi 28 $N$ trong $2\,\text{s}$. Bỏ qua lực cản. Tính quãng đường.
\loigiai{
Phân tích dữ kiện, chọn định luật phù hợp. $a=F/m=9,33m/s^2;s=\frac12at^2=18,66$ m.
}
\end{bt}
